\documentclass[11pt]{article}
\usepackage[T1]{fontenc}
\usepackage{lmodern,microtype}
\usepackage[margin=1in]{geometry}
\usepackage{amsmath,amssymb,amsthm,mathtools}
\usepackage{booktabs,longtable,array,multicol}
\usepackage{xurl}
\usepackage[hidelinks]{hyperref}
\hypersetup{pdftitle={Finite real counterexamples and spectral obstructions for q-permanent monotonicity},pdfauthor={Dev D. Goyal}}
\newtheorem{theorem}{Theorem}[section]
\newtheorem{lemma}[theorem]{Lemma}
\newtheorem{proposition}[theorem]{Proposition}
\newtheorem{corollary}[theorem]{Corollary}
\theoremstyle{remark}\newtheorem{remark}[theorem]{Remark}
\newcommand{\per}{\operatorname{per}}
\newcommand{\inv}{\operatorname{inv}}
\newcommand{\R}{\mathbb R}
\newcommand{\C}{\mathbb C}
\newcommand{\Z}{\mathbb Z}
\newcommand{\F}{\mathrm F}
\newcommand{\NN}[2]{\mathcal N_{#1}(#2)}
\newcommand{\WitnessN}{896}

\title{Finite real counterexamples and spectral obstructions\protect\\
for $q$-permanent monotonicity}
\author{Dev D. Goyal}
\date{1 November 2026}
\begin{document}
\maketitle
\begin{abstract}
For an ordered real matrix, write
$P_q(A)=\sum_{\sigma}q^{\inv(\sigma)}\prod_i a_{i,\sigma(i)}$.
We give an explicit integer Gram matrix of order $\WitnessN$ and rank three
whose derivative at $q=1$ is strictly negative. Its certificate uses a
four-row symmetry and exact integer polynomial arithmetic; all input
vectors are listed. An explicit scalar multiple plus the identity is an
integer positive-definite counterexample, with a fully specified rational
decreasing pair in $(0,1)$, and identity padding gives real
positive-definite counterexamples in every order $N\ge\WitnessN$. This
answers the real-symmetric case of the Bapat--Lal monotonicity
conjecture on $[-1,1]$ negatively, and refutes da Fonseca's extended
conjectures inside the original interval. We also derive the complete limiting endpoint
functional for bounded fixed-rank perturbations $I+(t/N)VV^T$.
For nonzero limiting covariance, below $t\lambda_{\max}(\Sigma)=1$ the
endpoint derivative is eventually positive for every ordering. Convergence
of the ordered empirical measure gives an explicit positive first-order limit.
An operator factorization of the complete endpoint score yields a stronger
spectral obstruction: a negative limit at scale $N^2\per(A_N)$ requires
at least three covariance eigenvalues strictly larger than $1/t$.
If $t\lambda_1(\Sigma)>1\ge t\lambda_3(\Sigma)$, the normalized
derivative instead has an explicit strictly positive lower bound, for
every ordering; this includes every supercritical real perturbation of
latent rank two.
For a convergent supercritical row cloud, uniformly random orders have a
fixed positive normalized derivative except on an exponentially small
fraction of labelled permutations.
Above the threshold, anisotropic spherical families yield correlation
counterexamples with off-diagonal entries $O(N^{-1})$, least eigenvalue
tending to one, and condition number tending to $227$ in latent rank four
or $20001$ in latent rank three. A Lean~4 development proves the
endpoint identity, the soundness of the integer certificate and the
reductions to rational positive-definite counterexamples in every order
$N\ge\WitnessN$; the finite integer comparison itself is checked by two
independent exact programs, and the analytic asymptotic proofs are
conventional.
\end{abstract}

\section{Results and context}

For $\sigma\in S_n$, let
\[
\inv(\sigma)=\#\{(i,j):i<j,\ \sigma(i)>\sigma(j)\},\qquad
P_q(A)=\sum_{\sigma\in S_n}q^{\inv(\sigma)}\prod_{i=1}^n a_{i,\sigma(i)}.
\]
The convention is $0^0=1$, and $D_q(A)=\frac{d}{dq}P_q(A)$.
The fixed order matters. A simultaneous relabelling by an arbitrary
permutation generally changes this polynomial.
At $q=-1,0,1$ it gives the determinant, the diagonal product, and the
permanent, respectively. Bapat and Lal's monotonicity conjecture asks
whether a non-diagonal positive-definite matrix has strictly increasing
$q$-permanent on $[-1,1]$ \cite{Bapat1992,BapatLal1994}.
Mitchell proved the equivalence with a positive-semidefinite formulation
with positive diagonal, and established rank-one and order-three cases
\cite{Mitchell2020}.

Da Fonseca conjectured strict monotonicity of the $q$-permanent of a
real symmetric positive-definite matrix on an interval with left
endpoint at or below $-1$ and right endpoint $+\infty$
\cite[Conjecture~1]{daFonseca2010}, and restated it with left endpoint
strictly below $-1$ \cite[Conjecture~2]{daFonseca2018}; both contain the
real-symmetric case of the Bapat--Lal assertion on $[-1,1]$.

Kitamura's public Lean proof supplies a complex Hermitian order-144
counterexample by a negative derivative at one \cite{Kitamura2026}.
The public problem record, consulted on 1 November 2026, explicitly
distinguishes that result from the real-symmetric case, which it lists
as unsettled \cite{OpenProblemsMI18}.
A companion paper gave real positive-definite counterexamples to da
Fonseca's conjectures in the extended range $q>1$
\cite{GoyalExtended2026}; it left the real-symmetric question on
$[-1,1]$ open. The present paper concerns that original interval,
with a negative derivative at its positive endpoint and strict decrease
at two parameters in $(0,1)$. The real-symmetric case of the Bapat--Lal
conjecture is therefore false, and da Fonseca's conjectures fail inside
the original interval as well as above it.
The mean-field theory below additionally explains an obstruction
to arbitrarily weak perturbations within the fixed-rank mechanism.

\begin{theorem}[A finite real endpoint counterexample]\label{thm:finite}
The ordered integer rows in Appendix~\ref{app:data}, interpreted by the
four-row rule in Section~\ref{sec:certificate}, define a Gram matrix
$B\in\Z^{\WitnessN\times\WitnessN}$ of rank exactly three, with positive
diagonal, such that
\[
 D_1(B)<0.
\]
The stronger exact certificate margin is recorded in
Table~\ref{tab:witness}. In particular there are a real positive-definite
matrix of this same order and rational $0<q_1<q_2<1$ with
$P_{q_2}(A)<P_{q_1}(A)$.
\end{theorem}

The positive-definite and rational-parameter conclusion is elementary:
the endpoint derivative is continuous in the entries and in $q$; adding
a sufficiently small positive diagonal preserves its strict negative
sign, and a small interval immediately to the left of one contains a
rational decreasing pair. Section~\ref{fr:section} strengthens this to a
specific integer positive-definite matrix and a specified pair.
It follows immediately that real positive-definite counterexamples
exist in every order $N\ge896$: use $A\oplus I_{N-896}$.
Every nonzero permutation term fixes the trailing identity coordinates,
and these coordinates introduce no inversions, so
$P_q(A\oplus I_{N-896})=P_q(A)$ for every real $q$.

\begin{theorem}[Bounded-condition asymptotic families]\label{thm:families}
There are explicit real algebraic positive-definite correlation matrices
$C_N$ in every sufficiently large order such that
\[
 \limsup_{N\to\infty}\frac{D_1(C_N)}{N(N-1)\per C_N}<-\frac1{40000},
 \qquad |(C_N)_{ij}|\le\frac{904}{N}\quad(i\ne j),
\]
and
\[
 \lambda_{\min}(C_N)\longrightarrow1,\qquad
 \kappa_2(C_N)\longrightarrow227.
\]
In particular $\kappa_2(C_N)<228$ eventually. Before diagonal
normalization, the matrices are $I_N+(904/N)V_NV_N^T$ with four real
columns in $V_N$.

A second four-column family has limiting condition number $256$ and
the larger normalized derivative margin $-1/1000$.

With three columns instead, another explicit family has off-diagonal
bound $60000/N$, limiting condition number $20001$, and normalized
endpoint derivative limit less than $-1/800$.
\end{theorem}

These are eventual-order conclusions, with no new explicit sufficient
order claimed. They are not assertions of convergence to the identity
in operator norm: finitely many eigenvalues remain separated from one.
They are also different families from the finite certificate.
The full proofs, including the algebraic grids and all constants, are
Theorems~\ref{mf:20002}, \ref{mf:257}, and~\ref{rf:228}.

\begin{theorem}[A strict fixed-rank obstruction]\label{thm:obstruction-intro}
Let uniformly bounded real rows of a fixed latent dimension satisfy
$N^{-1}V_N^TV_N\to\Sigma\ne0$. For fixed $t>0$, if
$t\lambda_{\max}(\Sigma)<1$, then
\[
 D_1\bigl(I_N+(t/N)V_NV_N^T\bigr)>0
\]
for all sufficiently large $N$, for every sequence of row orderings.
The same holds after correlation normalization. With convergence of the
ordered empirical measure, an explicit strictly positive limit of
$D_1/(N\per)$ is given in Theorem~\ref{mf:subcritical}.
\end{theorem}

Thus limiting condition number below two is an obstruction for this
bounded-row, fixed-rank scaling. This theorem does not extend its
conclusion to every positive-definite matrix of condition number below
two. Nor is the coherent-saddle transition automatically the transition
in the sign of $D_1$.
Indeed, Corollary~\ref{mf:product-limit} shows that a limiting row law
independent of position gives a strictly positive leading endpoint
derivative even above that threshold. The complete ordered functional
therefore distinguishes these positive configurations from the
counterexample orderings.

\begin{theorem}[A third-eigenvalue obstruction]\label{thm:third-intro}
Let $r\ge1$ be fixed, $\|v_{i,N}\|\le L$, and
$N^{-1}V_N^TV_N\to\Sigma$. List the covariance eigenvalues in
decreasing order, setting $\lambda_3(\Sigma)=0$ when $r<3$.
For fixed $t>0$, if
\[
 t\lambda_1(\Sigma)>1\qquad\hbox{and}\qquad
 t\lambda_3(\Sigma)\le1,
\]
then, for every sequence of row orderings,
\[
 \liminf_{N\to\infty}
 \frac{D_1(I_N+(t/N)V_NV_N^T)}
 {N(N-1)\per(I_N+(t/N)V_NV_N^T)}
 \ge\frac{(t\lambda_1(\Sigma)-1)^2}{64r\,t^2L^4}>0.
\]
The same normalized ratio is unchanged by correlation normalization.
\end{theorem}

Thus a large largest eigenvalue alone is insufficient for the negative
endpoint mechanism. For a strictly negative $N^2$-scale normalized
liminf, at least three covariance modes must exceed the
threshold $1/t$. The theorem allows equality for the third eigenvalue
and places no upper bound on the first two eigenvalues. In latent rank
two its explicit lower bound has denominator $128t^2L^4$; together with
Theorem~\ref{thm:obstruction-intro}, this gives eventual positivity at
every noncritical strength. Section~\ref{sp:section} proves the stronger
inertia statement behind this obstruction and treats its exact scope.

\paragraph{Contribution and verification scope.}
The Gaussian representation of a permanent is standard; applications
to spatial permutations and related asymptotics appear, for example,
in Fyodorov and Muirhead \cite{FyodorovMuirhead2021}. The new results here
are the finite real certificate, its symmetry reduction, the complete
ordered endpoint functional, the strict subcritical asymptotic, and
the stated bounded-condition counterexamples, together with the full-score
operator factorization, quantitative contraction gap, and third-eigenvalue
obstruction. Literature searches
before and after the derivations found no matching prior results in
the primary sources examined. This is a bounded novelty assessment,
not a certification of priority.
The ordinary permanent-on-top and subset-preserving $q$-permanent
inequalities are different statements; counterexamples to them are
not used as counterexamples to monotonicity.
A public manuscript of liqsweep claims permanent-on-top for real
positive-semidefinite matrices of rank at most two \cite{Liqsweep2026}.
As discussed in Section~\ref{sp:prior}, this has a qualitative endpoint
consequence after diagonal addition. We do not claim priority for that
rank-two qualitative sign. The proof of Theorem~\ref{thm:third-intro}
is independent of that manuscript and gives an explicit asymptotic
margin in arbitrary fixed latent dimension.
Corollary~\ref{ro:random-order} further shows that negative endpoint
orders occupy an exponentially small fraction of labelled permutations
of a convergent supercritical row cloud. This is a consequence of the
product-limit formula and bounded-difference concentration, not a claim
that the counterexample orders have disappeared.

\section{The endpoint identity}\label{sec:endpoint}

Let $v_1,\ldots,v_n\in\R^r$ be ordered rows, $B_{ij}=v_i\cdot v_j$,
and $\ell_i(z)=v_i\cdot z$. For a polynomial $p=\sum_\alpha c_\alpha z^\alpha$,
its Fischer norm is
\[
 \|p\|_{\F}^2=\sum_\alpha\alpha!|c_\alpha|^2.
\]
For an exterior-vector-valued polynomial, sum these norms over an
orthonormal exterior basis. Put
\[
 f=\prod_i\ell_i,\qquad
 G=\sum_{i<j}(v_i\wedge v_j)\prod_{a\ne i,j}\ell_a.
\]

\begin{lemma}[Gram and ordered wedge identities]\label{lem:endpoint}
One has
\[
 \per B=\|f\|_{\F}^2,
 \qquad
 2D_1(B)=\binom n2\|f\|_{\F}^2-\|G\|_{\F}^2.
\]
\end{lemma}
\begin{proof}
Expanding the product polynomial, a monomial exponent records the number
of rows assigned to each coordinate. Pairing two assignments with the
same exponent gives $\alpha!$ bijections of equal colors, exactly the
terms in the permanent expansion. This proves the first identity,
and the same argument applies to complementary row and column lists.

Expand $\|G\|_{\F}^2$ and use
\[
 \langle v_i\wedge v_j,v_k\wedge v_l\rangle
   =b_{ik}b_{jl}-b_{il}b_{jk}.
\]
For a fixed permutation monomial $\prod_i b_{i,\sigma(i)}$ and a fixed
row pair $i<j$, the contribution is $+1$ if
$\sigma(i)<\sigma(j)$ and $-1$ if $\sigma(i)>\sigma(j)$.
Summing over row pairs gives $\binom n2-2\inv(\sigma)$.
Thus
\[
 \|G\|_{\F}^2
 =\sum_\sigma\left(\binom n2-2\inv(\sigma)\right)
                         \prod_i b_{i,\sigma(i)},
\]
which is the second identity. This is the ordered-wedge formulation
used in \cite{Kitamura2026}; our real-rank-three coefficient and
polynomial versions are included in the formal development.
\end{proof}

For nonzero rows $f\ne0$, so its norm and the permanent are positive.
Consequently the dimensionless score
\[
 R(B)=\frac{\|G\|_{\F}^2}{\binom n2\per B}
\]
exceeds one exactly when $D_1(B)<0$, and
\begin{equation}\label{eq:normalized}
 \frac{D_1(B)}{n(n-1)\per B}=\frac{1-R(B)}4.
\end{equation}
Both sides are unchanged by scaling each row by an arbitrary nonzero
real number. This follows since every permutation monomial is
multiplied by the same positive product of squared row scales.

\section{A four-row integer certificate}\label{sec:certificate}

\subsection{The word and its compressed recurrence}

For a type $(a,b,c)$, define the two ordered pairs
\[
 L=((-a,b,c),(-a,-b,c)),\qquad
 R=((a,b,c),(a,-b,c)).
\]
Starting with an empty word, read the type list from inner to outer and
replace the current word $W$ by $L,W,R$. This completely specifies the
row order. The resulting word has four rows per type.

Write $X=x^2,Y=y^2,Z=z^2$ and
\begin{equation}\label{eq:orbitfactor}
 p=a^4X^2+b^4Y^2+c^4Z^2
       -2a^2b^2XY-2a^2c^2XZ-2b^2c^2YZ.
\end{equation}
The product polynomial is $f(x,y,z)=F(X,Y,Z)$. Its actual wedge has
the parity decomposition
\begin{align*}
 G_{01}&=yz\,U_{01}(X,Y,Z)+xz\,V_{01}(X,Y,Z),\\
 G_{02}&=U_{02}(X,Y,Z),\\
 G_{12}&=xy\,U_{12}(X,Y,Z)+V_{12}(X,Y,Z).
\end{align*}
Initially $F=1$ and the other five polynomials are zero. All arithmetic
below is over $\Z$. Subscripts denote partial derivatives in the
squared variables. One outer type updates the six polynomials by
\begin{align}
 F'&=pF,\label{eq:recF}\\
 U'_{01}&=pU_{01}+8ac\{b^2F+2b^2XF_X
                  +(a^2X+b^2Y-c^2Z)F_Y\},\label{eq:rec01}\\
 U'_{02}&=pU_{02}+8ac\{(a^2X-c^2Z)F
          +X(a^2X-b^2Y-c^2Z)F_X\notag\\
        &\hspace{39mm}+Z(a^2X+b^2Y-c^2Z)F_Z\},\label{eq:rec02}\\
 U'_{12}&=pU_{12}+8ac\{-b^2F
                 +(a^2X-b^2Y-c^2Z)F_Y-2b^2ZF_Z\},\label{eq:rec12}\\
 V'_{01}&=pV_{01}+8a^2bcF,\label{eq:recV01}\\
 V'_{12}&=pV_{12}+4bc(a^2X-b^2Y+c^2Z)F.\label{eq:recV12}
\end{align}

\begin{lemma}[Exact recurrence]\label{lem:recurrence}
Equations~\eqref{eq:recF}--\eqref{eq:recV12} give the ordinary product
and actual ordered wedge polynomials of the specified row word.
\end{lemma}
\begin{proof}
For concatenated words with pairs $(f,G)$ and $(h,H)$, separate wedge
pairs within each word and crossing between them. The cross sum is
$\nabla f\wedge\nabla h$, so concatenation gives
\[
 (f,G)\mathbin{\|}(h,H)=(fh,hG+fH+\nabla f\wedge\nabla h).
\]
For the negative and positive pairs their products are
$q_-=(cz-ax)^2-b^2y^2$ and $q_+=(cz+ax)^2-b^2y^2$;
their wedge vectors are $(2ab,0,2bc)$ and $(-2ab,0,2bc)$.
Applying concatenation twice gives
\[
 f'=q_-q_+f,
 \quad G'=pG+f(H+E)+\nabla f\wedge K,
\]
where direct differentiation gives
\begin{align*}
 H&=(8ab^2c\,yz,\ 8ac(a^2X-c^2Z),\ -8ab^2c\,xy),\\
 E&=(8a^2bc\,xz,\ 0,\ 4bc(a^2X-b^2Y+c^2Z)),\\
 K&=4ac\bigl(z(-a^2X-b^2Y+c^2Z),\ 2b^2xyz,\
                    x(a^2X-b^2Y-c^2Z)\bigr).
\end{align*}
Use $f_x=2xF_X,f_y=2yF_Y,f_z=2zF_Z$ and collect the five
indicated parity classes. This gives the displayed recurrences and
proves the assertion by induction.
\end{proof}

For $e\in\{0,1\}^3$ and $Q=\sum_\alpha Q_\alpha X^{\alpha_0}Y^{\alpha_1}Z^{\alpha_2}$,
put
\[
 \mathcal N_e(Q)=\sum_\alpha Q_\alpha^2
                   (2\alpha_0+e_0)!(2\alpha_1+e_1)!(2\alpha_2+e_2)!.
\]
Distinct physical exponent parities are Fischer-orthogonal. Hence
\begin{align}
 P:=\per B&=\NN{000}{F},\label{eq:certP}\\
 T_U&:=\NN{011}{U_{01}}+\NN{000}{U_{02}}+\NN{110}{U_{12}},\label{eq:certU}\\
 T:=\|G\|_{\F}^2&=T_U+\NN{101}{V_{01}}+\NN{000}{V_{12}}\ge T_U.
 \label{eq:certT}
\end{align}
The exact integer inequality $T_U>\binom n2P$ therefore suffices for
a negative endpoint derivative. The two correction terms need not be
evaluated for that conclusion; they are retained in the independent
full-wedge arithmetic check.

\subsection{The finite data and sign}

\begin{table}[htbp]
\centering
\caption{Exact finite witness. Ratios displayed as decimals are explanatory;
the strict rational inequality is checked on the full integers.}
\label{tab:witness}
\begin{tabular}{@{}ll@{}}
\toprule
Quantity & Value or exact bound\\
\midrule
Order / number of four-row types & $896$ / $224$\\
Rank of the integer Gram matrix & $3$\\
Largest absolute row coordinate & $32$\\
Largest Gram diagonal entry & $2490$\\
Nonzero row minor (one-based indices) & $(417,418,477)$, determinant $1012$\\
Number of decimal digits of $P$ & $4152$\\
Three-component ratio $T_U/(\binom{896}{2}P)$ & $1.0019230541058451\ldots$\\
Full-wedge ratio $T/(\binom{896}{2}P)$ & $1.0019250593839399\ldots$\\
Exact strict certificate & $1000T_U>1001\binom{896}{2}P$\\
Consequence for $D_1(B)/\per B$ & $<-5012/25$\\
Consequence for $D_1(B)/(896\cdot895\per B)$ & $<-1/4000$\\
Explicit positive-definite scaling $M=10^4n^3$ & $7193231360000$\\
Entry bound for $A=MB+I$: $H=2490M+1$ & $17911146086400001$\\
\bottomrule
\end{tabular}
\end{table}

The exact report stores $P,T_U,T$ and all five component norms as decimal
integers. The SHA-256 of the inner-to-outer input file is
\begin{center}\small\texttt{7220026a9fb5b2d9e62ef46b9b7196b2ff768b4841f71ccb45275ee445987751}.\end{center}


The data in Appendix~\ref{app:data} specify every row, without any
optimization routine or floating-point number in their definition.
Applying the integer recurrences and factorial sums gives the exact
strict inequalities in Table~\ref{tab:witness}. Lemma~\ref{lem:endpoint}
then proves the negative sign. Positive semidefiniteness follows from
$B=VV^T$. A nonzero three-by-three minor specified in the table proves
rank three; each row is nonzero, so every diagonal entry is positive.
This proves Theorem~\ref{thm:finite}.

The accompanying package also records the full integer numerators and
denominator of the rational pair specified in Section~\ref{fr:section}.
The common denominator has $16834$ decimal digits; its compact exact
formula is preferable in the mathematical text.

The exact engine compares integers, not the displayed decimal ratios.
A second, independent program uses neither the four-row symmetry nor
the recurrences above. It expands the ordinary product and ordered-wedge
recurrences
\[
 f'=f\ell,\qquad S'_a=S_a\ell+v_af,\qquad
 G'_{ab}=G_{ab}\ell+v_bS_a-v_aS_b
\]
row by row in the actual order, where $\ell=v\cdot z$ is the new row
form and $S_a=\partial f/\partial z_a$, and evaluates $P$ and the
Fischer norms of every exponent-parity class of the three components
of $G$ modulo $500$ primes between $2^{30}$ and $2^{31}$. Every permutation
monomial has absolute value at most $\prod_ib_{ii}$, so
$P\le n!\prod_ib_{ii}$ and each parity-class norm is at most
$\binom n2n!\prod_ib_{ii}$. The product of the primes exceeds this
bound, which has $15471$ bits. The Chinese remainder theorem therefore
determines $P$, $T_U$ and $T$ exactly from the residues, and all three
agree with the integer report. The only nonzero parity classes are the
five in \eqref{eq:certP}--\eqref{eq:certT}. This determination is
independent of Lemma~\ref{lem:recurrence}; it relies only on
Lemma~\ref{lem:endpoint} and exact arithmetic. The precise Lean proof
boundary is given in Section~\ref{sec:lean}.

% Input-ready conventional-proof component.
% Requires amsmath, amssymb, amsthm and theorem/lemma/corollary environments.
% This component is independent of the version of the finite witness data.
% All labels in this file have the prefix fr:.
\providecommand{\per}{\operatorname{per}}

\section{An explicit integer positive-definite counterexample from a finite certificate}
\label{fr:section}

This section converts a finite positive-semidefinite endpoint certificate
with a specified margin into an integer positive-definite matrix and
two specified rational parameters in \((0,1)\). The estimates below are
conventional proofs; their explicit numerical regularization is separate
from the scope of the Lean certificates reported elsewhere.

For an ordered \(k\)-by-\(k\) positive-semidefinite matrix \(C\), write
\(D_1(C)=P'_1(C)\), \(c_k=\binom{k}{2}\), and
\[
 T(C)=\left\|\sum_{i<j}(v_i\wedge v_j)
                    \prod_{a\ne i,j}(v_a\cdot z)\right\|_{\mathrm F}^2,
 \qquad C=(v_i\cdot v_j)_{i,j}.
\]
Thus \(2D_1(C)=c_k\per C-T(C)\). The empty permanent is one.
For \(I\subseteq[n]\), let \(D_{\rm emb}(B[I])\) denote the derivative
obtained by allowing permutations on \(I\) and fixing every index in
\([n]\setminus I\), with inversions counted in the original order on
\([n]\).

\begin{lemma}[A uniform bound for embedded principal derivatives]
\label{fr:embedded-bound}
Let \(B\) be a real positive-semidefinite matrix of order \(n\ge2\),
let \(I\subseteq[n]\), and put \(k=|I|\), \(j=n-k\). Then
\begin{equation}\label{fr:embedded-estimate}
 |D_{\rm emb}(B[I])|
 \le \left\{\frac{c_k+c_k^2}{2}+\frac{jk^2}{2}\right\}\per B[I]
 \le n^4\per B[I].
\end{equation}
\end{lemma}

\begin{proof}
For \(C=B[I]\), define the permanental cofactor matrix
\[
 F_{ab}(C)=c_{ab}\per C(a\mid b),
\]
where \(C(a\mid b)\) deletes row \(a\) and column \(b\).
The matrix \((\per C(a\mid b))_{a,b}\) is a Gram matrix of the
products of all but one row form, in the Fischer inner product.
The Schur product theorem therefore shows that \(F(C)\) is positive
semidefinite. The singleton case of Lieb's permanent inequality
\cite{Lieb1966} gives
\(F_{aa}(C)\le\per C\), and hence
\begin{equation}\label{fr:F-bound}
 |F_{ab}(C)|\le\per C.
\end{equation}
For each fixed index \(u\notin I\), write
\[
 \operatorname{cut}_{u}F(C)
   =\sum_{\substack{a,b\in I\\a<u<b}}F_{ab}(C),
\]
with the labels on \(I\) inherited from \([n]\). Inserting the fixed
index \(u\) contributes twice the number of active indices mapped from
its left to its right. Summing over permutations gives the exact identity
\begin{equation}\label{fr:padding}
 D_{\rm emb}(C)=D_1(C)+2\sum_{u\notin I}\operatorname{cut}_{u}F(C).
\end{equation}
There are at most \(k^2/4\) pairs crossing each cut. Thus
\eqref{fr:F-bound} bounds the second term in \eqref{fr:padding} in
absolute value by \(jk^2\per C/2\).

Each of the \(c_k\) summands in the exterior polynomial defining
\(T(C)\) has squared Fischer norm
\[
 (c_{aa}c_{bb}-c_{ab}^2)\per C[\{a,b\}^c]
 \le c_{aa}c_{bb}\per C[\{a,b\}^c]\le\per C.
\]
The final inequality follows by applying the singleton form of Lieb's
inequality twice. The triangle inequality, or Cauchy--Schwarz, now gives
\(T(C)\le c_k^2\per C\). Therefore
\[
 |D_1(C)|\le\frac{c_k+c_k^2}{2}\per C.
\]
The same statement holds for \(k=0,1\), when both sides involving the
derivative vanish. These estimates prove the first inequality in
\eqref{fr:embedded-estimate}. For \(k\ge1\),
\((c_k+c_k^2)/2\le k^4/2\), while
\(jk^2/2\le n^3/2\le n^4/2\). This proves the second inequality;
the case \(k=0\) is immediate.
\end{proof}

\begin{theorem}[Explicit integer regularization]
\label{fr:integer-regularization}
Let \(n\ge6\), and let \(B\) be an integer real positive-semidefinite
matrix of order \(n\) whose diagonal entries are at least one. Suppose
that
\begin{equation}\label{fr:input-margin}
 T(B)>\frac{1001}{1000}\binom n2\per B.
\end{equation}
Set
\begin{equation}\label{fr:explicit-matrix}
 M=10^4n^3,\qquad A=MB+I_n,\qquad
 H_0=\max_i b_{ii},\qquad H=MH_0+1.
\end{equation}
Then \(A\) is an integer positive-definite matrix, its entries have
absolute value at most \(H\), and
\begin{equation}\label{fr:retained-margin}
 D_1(A)<-M^n\frac{n(n-5)}{20000}\per B<0.
\end{equation}
Its least eigenvalue is at least one, and is exactly one when \(B\)
is singular. If \(B\) has rank \(r<n\), then precisely \(n-r\)
eigenvalues of \(A\) equal one.
\end{theorem}

\begin{proof}
Put \(\eta=M^{-1}\). Expansion in the diagonal additions gives
\begin{equation}\label{fr:diagonal-expansion}
 D_1(B+\eta I_n)
 =D_1(B)+\sum_{\varnothing\ne J\subseteq[n]}
       \eta^{|J|}D_{\rm emb}(B[J^c]).
\end{equation}
The diagonal assumptions and the iterated singleton form of Lieb's
inequality imply
\(\per B[J^c]\le\per B\) for every \(J\). They also imply
\(\per B\ge\prod_i b_{ii}\ge1\). By
Lemma~\ref{fr:embedded-bound}, the sum in
\eqref{fr:diagonal-expansion} has absolute value at most
\[
 n^4\per B\bigl((1+\eta)^n-1\bigr).
\]
Since \(n\eta\le1/2\), the inequalities
\(\binom nj\le n^j\) and the geometric-series formula give
\[
 (1+\eta)^n-1
 \le\sum_{j=1}^\infty(n\eta)^j
 =\frac{n\eta}{1-n\eta}\le2n\eta.
\]
The perturbation in the derivative is consequently at most
\(2n^5\eta\per B=n^2\per B/5000\). On the other hand,
\eqref{fr:input-margin} and the endpoint identity imply
\[
 D_1(B)<-\frac{n(n-1)}{4000}\per B.
\]
Combining the last two estimates yields
\[
 D_1(B+\eta I_n)
 <-\left(\frac{n(n-1)}{4000}-\frac{n^2}{5000}\right)\per B
 =-\frac{n(n-5)}{20000}\per B<0.
\]
The derivative is homogeneous of degree \(n\) in the matrix entries,
so multiplication by \(M^n\) gives \eqref{fr:retained-margin}.
The remaining assertions follow from \(A=MB+I_n\), from the spectral
theorem, and from
\(|b_{ij}|\le\sqrt{b_{ii}b_{jj}}\le H_0\).
\end{proof}

A certificate may establish \eqref{fr:input-margin} by computing any
explicit lower bound \(U\le T(B)\) and checking the integer inequality
\(1000U>1001\binom n2\per B\). In particular, a sum of selected mutually
orthogonal parity-component norms is sufficient. The values of \(n\)
and \(H_0\) are read from the particular frozen row file; none of the
preceding estimates depends on a choice between witness versions.

\begin{corollary}[Specified rational points of strict decrease]
\label{fr:rational-interior}
Under the hypotheses of Theorem~\ref{fr:integer-regularization}, define
the positive integer
\begin{equation}\label{fr:Q}
 Q=4n!\,n^4H^n,
 \qquad q_1=1-\frac2Q,\qquad q_2=1-\frac1Q.
\end{equation}
Then \(0<q_1<q_2<1\), and
\begin{equation}\label{fr:strict-interior}
 P'_{q}(A)\le-\frac12\quad(q_1\le q\le1),
 \qquad P_{q_1}(A)>P_{q_2}(A).
\end{equation}
Thus the integer matrix \(A\) and the two rational parameters in
\eqref{fr:Q} form a fully specified positive-definite counterexample.
\end{corollary}

\begin{proof}
The polynomial \(P_q(A)\) has integer coefficients. Since
\(D_1(A)<0\), it follows that \(D_1(A)\le-1\).
For \(0\le q\le1\), termwise differentiation of the permutation
expansion and \(\operatorname{inv}(\sigma)\le\binom n2\le n^2\)
give
\[
 |P''_q(A)|\le n!\,n^4H^n=\frac Q4.
\]
Terms with fewer than two inversions contribute zero, so this bound
also holds at \(q=0\). For \(q\in[1-2/Q,1]\), the fundamental
theorem of calculus gives
\[
 P'_q(A)\le P'_1(A)+\frac Q4(1-q)\le-1+\frac12=-\frac12.
\]
Since \(Q>2\), the indicated interval lies in \((0,1]\), and
integration between \(q_1\) and \(q_2\) proves the strict inequality.
\end{proof}

% Input after finite_regularization.tex and before mean_field.tex.
% Conventional proof; labels use co:. No later-section dependencies.

\section{A one-parameter reduction for coherent ordering}
\label{co:section}

The four-row symmetry also permits an exact ordering theorem for a
single coherent evaluation. This theorem concerns the actual exterior
polynomial, including its correction terms.

Take \(m\) triples \((a_j,b_j,c_j)\) of nonnegative real numbers,
with \(a_j^2+b_j^2>0\) for each triple, and assume that at least one
triple has \(a_jb_jc_j>0\). An order of the triples defines a word
of \(4m\) rows by starting with the empty word and successively
prepending \((-a_j,b_j,c_j),(-a_j,-b_j,c_j)\) and appending
\((a_j,b_j,c_j),(a_j,-b_j,c_j)\). Indices \(j=0,\ldots,m-1\)
refer to this inner-to-outer order. For its rows \(v_i\), set
\[
 f(z)=\prod_i(v_i\cdot z),\qquad
 G_{ab}(z)=\sum_{i<j}(v_{ia}v_{jb}-v_{ib}v_{ja})
                         \prod_{h\ne i,j}(v_h\cdot z),
 \quad 0\le a<b\le2.
\]

Fix \(s>0\), put \(z_s=(1,is,0)\), and define
\begin{align*}
 S_j&=a_j^2+s^2b_j^2,&
 t_j&=\frac{a_jc_j}{S_j},&
 d_j&=\frac{a_j^2-s^2b_j^2}{S_j},&
 K_s&=\sum_j\frac{b_jc_j}{S_j},\\
 A&=\sum_j(2j+1)t_j,&
 B&=\sum_jt_j\left(d_j+2\sum_{\ell<j}d_\ell\right).
\end{align*}
The quantities \(A,B\) depend on the order; \(K_s\) does not.

\begin{proposition}[Exact coherent evaluation]\label{co:evaluation}
One has \(f(z_s)=\prod_jS_j^2>0\), and
\begin{equation}\label{co:components}
 \frac{G_{01}(z_s)}{f(z_s)}=0,\qquad
 \frac{G_{02}(z_s)}{f(z_s)}=4(A+B),\qquad
 \frac{G_{12}(z_s)}{f(z_s)}=4K_s+\frac{4i}{s}(B-A).
\end{equation}
Consequently
\begin{equation}\label{co:score}
 \mathcal E_s:=
 \frac{\sum_{a<b}|G_{ab}(z_s)|^2}{|f(z_s)|^2}
 =16\left((A+B)^2+\frac{(B-A)^2}{s^2}+K_s^2\right).
\end{equation}
At the unit point \(z_s/\sqrt{1+s^2}\), the corresponding ratio is
\((1+s^2)^2\mathcal E_s\).
\end{proposition}
\begin{proof}
The polynomial identities in the four-row recurrence hold for real
inputs as well as integers. Evaluate the squared variables at
\((X,Y,Z)=(1,-s^2,0)\). For the orbit factor \(p_j\),
\[
 p_j=S_j^2,\qquad
 \frac{(p_j)_X-s^2(p_j)_Y}{p_j}=2,\qquad
 \frac{(p_j)_X+s^2(p_j)_Y}{p_j}=2d_j.
\]
Thus, for the product \(F\) of the first \(j\) orbit factors,
\[
 \frac{F_X-s^2F_Y}{F}=2j,\qquad
 \frac{F_X+s^2F_Y}{F}=2\sum_{\ell<j}d_\ell.
\]
Substituting in \eqref{eq:rec02}--\eqref{eq:rec12} gives the
increments of the normalized ratios
\begin{align*}
 \Delta\frac{U_{02}-s^2U_{12}}F&=8(2j+1)t_j,\\
 \Delta\frac{U_{02}+s^2U_{12}}F
 &=8t_j\left(d_j+2\sum_{\ell<j}d_\ell\right).
\end{align*}
Equation~\eqref{eq:recV12} gives
\(\Delta(V_{12}/F)=4b_jc_j/S_j\).
Summing from the zero initial exterior polynomials and using
\(G_{02}=U_{02}\), \(G_{12}=xyU_{12}+V_{12}\), and
\(G_{01}=yzU_{01}+xzV_{01}\) proves \eqref{co:components}.
Its real and imaginary parts give \eqref{co:score}. Finally \(G\)
has degree two less than \(f\), which proves the unit-point scaling.
\end{proof}

\begin{theorem}[One-parameter sorting]\label{co:sorting}
For every fixed \(s>0\), an ordering maximizing \(\mathcal E_s\)
is obtained by sorting the triples in nondecreasing order of
\begin{equation}\label{co:keys}
 \frac{a_jc_j}{a_j^2+\kappa b_j^2}
\end{equation}
for some \(\kappa>0\), with suitable ordering among ties. The
family of sorts is independent of \(s\); its maximizing parameter
may depend on \(s\).
\end{theorem}
\begin{proof}
Since \(t_j\ge0\) and \(|d_j|\le1\), one has \(|B|\le A\).
The triple with all entries positive has \(t_j>0\) and
\(|d_j|<1\), so its contribution makes \(|B|<A\) in every order.
The order-independent term \(K_s^2\) can be dropped when maximizing.
Write
\[
 Q(A,B)=(A+B)^2+s^{-2}(B-A)^2
       =D(A^2+B^2)+2CAB,
 \qquad D=1+s^{-2},\quad C=1-s^{-2}.
\]
Here \(D>|C|\), so \(Q\) is positive definite.

Choose an order maximizing \(Q\), and denote its pair by \((A,B)\).
It also maximizes the linear functional \(A'+\lambda B'\), where
\begin{equation}\label{co:lambda}
 \lambda=\frac{CA+DB}{DA+CB}\in(-1,1).
\end{equation}
Indeed, \(DA+CB>0\), and a strict increase in this functional
would strictly increase \(Q\) by its quadratic expansion.
Moreover
\[
 (DA+CB)\mathbin\pm(CA+DB)=(D\pm C)(A\pm B)>0,
\]
which proves the strict parameter bounds in \eqref{co:lambda}.

If adjacent types \(i,j\) occur in that order, swapping them changes
\(A+\lambda B\) by
\[
 2\bigl((1+\lambda d_j)t_i-(1+\lambda d_i)t_j\bigr).
\]
Both factors \(1+\lambda d_j\) are positive. A strict inversion
in the keys \(t_j/(1+\lambda d_j)\) would therefore improve the
functional, contradicting maximality. The order is sorted in those
keys, and
\[
 \frac{t_j}{1+\lambda d_j}
 =\frac1{1+\lambda}
   \frac{a_jc_j}
        {a_j^2+s^2\frac{1-\lambda}{1+\lambda}b_j^2}.
\]
Taking \(\kappa=s^2(1-\lambda)/(1+\lambda)>0\) proves the theorem.
\end{proof}

\begin{corollary}[A finite sweep, including ties]\label{co:sweep}
At most \(\binom m2+1\) interval sorts suffice to attain the global
maximum in Theorem~\ref{co:sorting}. If the row coordinates and \(s^2\)
are rational, the sweep and the score comparisons can be carried out
with exact rational arithmetic.
\end{corollary}
\begin{proof}
Put \(u_j=a_jc_j\). Two keys that are not identical as functions
of \(\kappa\) cross at most once for \(\kappa>0\), at
\[
 \kappa=\frac{u_j a_i^2-u_i a_j^2}
                  {u_i b_j^2-u_j b_i^2},
\]
provided the denominator is nonzero and the quotient is positive.
These crossings partition \((0,\infty)\) into at most
\(\binom m2+1\) open intervals. Sort once at any parameter in
each interval.

To check that crossing ties cannot hide a larger value, let \(X=(A,B)\)
be a global maximizer and use its gradient functional from
\eqref{co:lambda}. Any other ordering \(Y\) tied for that functional
must have \(Y=X\): otherwise the quadratic expansion gives
\(Q(Y)=Q(X)+Q(Y-X)>Q(X)\).
There are finitely many orders. A sufficiently small one-sided
perturbation of the parameter cannot promote an order whose supporting
value was strictly smaller, so it selects a tied maximizer with the
same \((A,B)\). This covers simultaneous crossings as well.
Keys identical for all parameters have zero adjacent-swap increment
in both \(A\) and \(B\), so arbitrary tie-breaking within those
classes does not change the score. All crossing parameters, keys, and
values of \(Q\) are rational under the stated assumptions.
\end{proof}

Individual zero values of \(a_j\), \(b_j\), or \(c_j\) are permitted;
the assumption of one strictly positive triple ensures \(A>|B|\) and
keeps the supporting parameter interior. If all \(t_j=0\), then
\(A=B=0\) and \(\mathcal E_s=16K_s^2\) is order-independent; this
trivial case does not need a supporting parameter. Boundary-only
systems with \(A>0\) but no strictly positive triple are outside the
sorting theorem as stated.

The optimization proved here is at one specified complex point.
A larger value of \(\mathcal E_s\) does not by itself imply a larger
coefficient-square Fischer norm or a negative endpoint derivative.
The finite signs elsewhere in this paper are certified separately by
their exact integer Fischer-norm computations.

% Input-ready conventional-proof component.
% Requires amsmath, amssymb, amsthm and theorem/lemma/proposition/corollary environments.
% Labels and equation labels in this file have the prefix mf:.
\providecommand{\per}{\operatorname{per}}

\section{Finite-rank mean-field perturbations}\label{mf:setup}

The next results give an asymptotic criterion for a negative endpoint
derivative and a strict obstruction in the subcritical regime. They are
conventional mathematical proofs. Their measure and limiting arguments
are not asserted to form a completed Lean theorem.

Fix a finite integer \(r\ge1\), a real number \(t>0\), and ordered vectors
\(v_{i,N}\in\mathbb R^r\), \(1\le i\le N\), with
\(\|v_{i,N}\|\le L\) for a constant \(L<\infty\) independent of \(N\).
Write \(V_N\) for the matrix with these rows and set
\begin{equation}\label{mf:matrices}
 B_N=V_NV_N^T,\qquad
 A_N=I_N+\frac tN B_N,\qquad
 \Sigma_N=\frac1N\sum_{i=1}^N v_{i,N}v_{i,N}^T.
\end{equation}
Every \(A_N\) is real positive definite. Suppose initially that the ordered
empirical measures satisfy
\begin{equation}\label{mf:empirical}
 \mu_N=\frac1N\sum_{i=1}^N\delta_{(i/N,v_{i,N})}
 \ \Longrightarrow\ \mu
\end{equation}
weakly on \([0,1]\times\overline B(0,L)\). The first marginal of \(\mu\)
is Lebesgue measure, and
\[
 \Sigma=\int vv^T\,d\mu(x,v)=\lim_N\Sigma_N.
\]
The order coordinate \(x\) is retained throughout.

For an ordered matrix \(A\), let \(D_1(A)=P'_1(A)\), with the
inversion-weighted \(q\)-permanent as defined in the main text.
For Gram rows \(w_i\), put
\[
 f(z)=\prod_i(w_i\cdot z),\qquad
 G(z)=\sum_{i<j}(w_i\wedge w_j)
                    \prod_{a\ne i,j}(w_a\cdot z),\qquad
 T(A)=\|G\|_{\mathrm F}^2.
\]
Here the Fischer norm is
\(\|\sum_\alpha c_\alpha z^\alpha\|_{\mathrm F}^2
=\sum_\alpha\alpha!|c_\alpha|^2\), summed over exterior components.
The identities
\begin{equation}\label{mf:endpoint}
 \per A=\|f\|_{\mathrm F}^2,\qquad
 2D_1(A)=\binom N2\per A-T(A)
\end{equation}
follow by expanding the Fischer inner products: the coefficient of a
permutation monomial in \(T(A)\) is
\(\binom N2-2\operatorname{inv}(\sigma)\).
The ordered wedge identity is also used in Kitamura's public proof
package \cite{Kitamura2026}.
Consequently, for positive permanent,
\begin{equation}\label{mf:ratio}
 \frac{D_1(A)}{N(N-1)\per A}
 =\frac14\left(1-\frac{T(A)}{\binom N2\per A}\right).
\end{equation}

\subsection{The general endpoint criterion}

For \(h\in\mathbb C^r\), define
\begin{align}
 \Phi(h)&=\int\log(1+t|v\cdot h|^2)\,d\mu(x,v)-\|h\|^2,
       &K&=\operatorname*{arg\,max}_{h\in\mathbb C^r}\Phi(h),\label{mf:potential}\\
 \ell_h(v)&=v\cdot h,
       &p_h(v)&=\frac{t|\ell_h(v)|^2}{1+t|\ell_h(v)|^2},\label{mf:p}\\
 q_h(v)&=\frac{t\overline{\ell_h(v)}\,v}{1+t|\ell_h(v)|^2},
       &m_h&=\int p_h(v)\,d\mu(x,v),\label{mf:q}\\
 Q_h(s)&=\int_{\{x<s\}}q_h(v)\,d\mu(x,v),
       &Q_h&=\int q_h(v)\,d\mu(x,v).\label{mf:prefix}
\end{align}
The vector-valued function \(q_h\) is continuous, including at
\(\ell_h(v)=0\). Put
\begin{equation}\label{mf:ordered-wedge}
 W_h=\iint_{\{x<y\}}q_h(v)\wedge q_h(w)\,
                       d\mu(x,v)\,d\mu(y,w)
\end{equation}
and
\begin{equation}\label{mf:score}
 \mathcal R(h)=2\|W_h\|^2+
 2\int(1-p_h(v))\|Q_h-2Q_h(x)\|^2\,d\mu(x,v)
 +(1-m_h)^2.
\end{equation}

\begin{theorem}[Full mean-field endpoint criterion]\label{mf:criterion}
Under \eqref{mf:empirical}, \(K\) is nonempty and compact, and
\begin{align}
 \min_{h\in K}\mathcal R(h)
 &\le\liminf_{N\to\infty}\frac{T(A_N)}{\binom N2\per A_N}\notag\\
 &\le\limsup_{N\to\infty}\frac{T(A_N)}{\binom N2\per A_N}
 \le\max_{h\in K}\mathcal R(h).\label{mf:criterion-bounds}
\end{align}
If \(\min_K\mathcal R>1\), then \(D_1(A_N)<0\) for all sufficiently
large \(N\). If \(\mathcal R=\mathcal R_*\) on \(K\), then
\begin{equation}\label{mf:criterion-limit}
 \frac{D_1(A_N)}{N(N-1)\per A_N}\longrightarrow
                  \frac{1-\mathcal R_*}{4}.
\end{equation}
At every \(h\in K\),
\begin{equation}\label{mf:stationarity}
 Q_h=\overline h,\qquad m_h=\|h\|^2<1.
\end{equation}
\end{theorem}

Both \(\Phi\) and \(\mathcal R\) are invariant under a global phase of
\(h\), and under conjugation because the rows are real.
Thus two conjugate phase circles of maxima give a single limiting score.
A simpler sufficient criterion is
\(\min_{h\in K}2\|W_h\|^2>1\). The other two terms in
\eqref{mf:score} are nonnegative.

\begin{corollary}[Product ordered limits]\label{mf:product-limit}
Suppose that the ordered empirical limit is
\(\mu(dx,dv)=dx\,\nu(dv)\), so that its row law does not depend on
the position coordinate. Then, for every \(h\in K\),
\begin{equation}\label{mf:product-score}
 \mathcal R(h)=1-\frac{m_h(4-m_h)}3.
\end{equation}
If \(t\lambda_{\max}(\Sigma)>1\), then
\begin{equation}\label{mf:product-positive}
 \liminf_{N\to\infty}
 \frac{D_1(A_N)}{N(N-1)\per A_N}
 \ge\min_{h\in K}\frac{m_h(4-m_h)}{12}>0.
\end{equation}
In particular the endpoint derivative is eventually positive.
\end{corollary}
\begin{proof}
The product form of \(\mu\) gives \(Q_h(s)=sQ_h\) and
\(W_h=\tfrac12 Q_h\wedge Q_h=0\). At a maximizer,
\eqref{mf:stationarity} gives \(\|Q_h\|^2=m_h\). Since \(p_h\)
also depends only on the row, the middle term of \eqref{mf:score} is
\[
 2(1-m_h)\|Q_h\|^2\int_0^1(1-2s)^2\,ds
 =\frac23m_h(1-m_h).
\]
Adding \((1-m_h)^2\) proves \eqref{mf:product-score}.

Let \(u\) be a real unit top eigenvector of \(\Sigma\). The elementary
bound \(\log(1+s)\ge s-s^2/2\) for \(s\ge0\) and the row bound
\(\|v\|\le L\) yield
\[
 \Phi(\epsilon u)
 \ge\epsilon^2\bigl(t\lambda_{\max}(\Sigma)-1\bigr)
       -\frac{t^2L^4}{2}\epsilon^4>0
\]
for sufficiently small \(\epsilon>0\). Since \(\Phi(0)=0\), zero
does not belong to \(K\). Compactness of \(K\) and
\(m_h=\|h\|^2<1\) therefore imply
\(\min_K m_h(4-m_h)>0\).
Combining \eqref{mf:product-score}, \eqref{mf:criterion-bounds}, and
\eqref{mf:ratio} proves \eqref{mf:product-positive}.
\end{proof}

The hypothesis is a product form of the ordered empirical limit itself;
no probabilistic assertion about finite random row orders is used.

\subsection{A strict obstruction below the covariance threshold}

\begin{theorem}[Subcritical endpoint increase]\label{mf:subcritical}
Suppose \(\Sigma\ne0\) and \(t\lambda_{\max}(\Sigma)<1\).
Define
\[
 R_0=(I_r-t\Sigma)^{-1},\qquad
 M_s=\int_{\{x<s\}}vv^T\,d\mu(x,v).
\]
Then
\begin{equation}\label{mf:permanent-limit}
 \per A_N\longrightarrow\det(I_r-t\Sigma)^{-1}
\end{equation}
and
\begin{equation}\label{mf:subcritical-limit}
 \frac{D_1(A_N)}{N\per A_N}\longrightarrow
 2t^2\int_0^1\operatorname{tr}
       \bigl(M_sR_0(\Sigma-M_s)\bigr)\,ds>0.
\end{equation}
In particular \(D_1(A_N)>0\) for all sufficiently large \(N\).
\end{theorem}

\begin{corollary}[Arbitrary changing orders]\label{mf:arbitrary-orders}
The ordered weak-convergence assumption can be dropped from the strict
sign conclusion. If only \(\Sigma_N\to\Sigma\ne0\) is assumed, with the
same fixed rank and row bound, and \(t\lambda_{\max}(\Sigma)<1\), then
\begin{equation}\label{mf:uniform-positive}
 \liminf_N\frac{D_1(A_N)}{N\per A_N}
 \ge\frac{t^2(\operatorname{tr}\Sigma^2)^2}{8L^4}>0.
\end{equation}
More generally, eventual bounds
\(t\lambda_{\max}(\Sigma_N)\le1-\varepsilon\) and
\(\operatorname{tr}(\Sigma_N^2)\ge s>0\), for fixed
\(\varepsilon,s>0\), imply the lower bound \(t^2s^2/(8L^4)\).
\end{corollary}

For \(N>r\), the least eigenvalue of \(A_N\) is one, while its largest
eigenvalue is \(1+t\lambda_{\max}(\Sigma_N)\).
Thus Theorem~\ref{mf:subcritical} gives a strict obstruction for this
entire bounded-row, fixed-rank mean-field construction below limiting
condition number two. It is not a theorem for every positive-definite
matrix of condition number below two, and does not include the critical
boundary \(t\lambda_{\max}(\Sigma)=1\).

\section{Proof of the full endpoint criterion}\label{mf:criterion-proof}

\subsection{Gaussian representation and concentration}

Gaussian representations of permanents also appear in the analysis of
spatial random permutations \cite{FyodorovMuirhead2021}.
We give the specific identities and limiting arguments needed here.
Let \(Z\) be a standard complex Gaussian vector in \(\mathbb C^r\), with
density \(\pi^{-r}e^{-\|z\|^2}\). Expanding by principal subsets and using
the Fischer identity gives
\begin{equation}\label{mf:gaussian}
 \per A_N=
 \mathbb E\prod_{i=1}^N\left(1+\frac tN|v_i\cdot Z|^2\right).
\end{equation}
After \(z=\sqrt N\,h\), the corresponding probability measure is
\begin{equation}\label{mf:tilted}
 \Pi_N(dh)=
 \frac{e^{N\Phi_N(h)}\,dh}{\int e^{N\Phi_N(u)}\,du},\qquad
 \Phi_N(h)=\frac1N\sum_i\log(1+t|v_i\cdot h|^2)-\|h\|^2.
\end{equation}
The functions \(\Phi_N\) converge to \(\Phi\) locally uniformly.
Indeed, weak convergence and uniform continuity on compact products give
uniform convergence on a finite \(h\)-net, and then on each compact set.
Uniformly in \(N\),
\begin{equation}\label{mf:coercive}
 \Phi_N(h)\le\log(1+tL^2\|h\|^2)-\|h\|^2.
\end{equation}
The right side is coercive and is at most \(-\|h\|^2/2\) outside a
fixed ball. This proves the compactness and nonemptiness of \(K\).

The measures \(\Pi_N\) concentrate in every neighborhood of \(K\).
To see this, take a large ball containing \(K\) and the region where
the upper bound in \eqref{mf:coercive} may be nonnegative.
On the compact complement of a neighborhood of \(K\), \(\Phi\) is
separated from its maximum by a strictly positive amount.
Local uniform convergence preserves a smaller gap for \(\Phi_N\).
A small ball around one maximizer gives a denominator bounded below by
its fixed volume times \(e^{N(\max\Phi-\epsilon)}\).
The outer tail is controlled uniformly by
\(e^{-N\|h\|^2/2}\). Choosing \(\epsilon\) smaller than the compact gap
proves the assertion. No nondegeneracy of the Hessian is needed.

Differentiating \(\Phi\) in the real and imaginary coordinates is
legitimate by the uniform row bound. Its stationary equation is
\(Q_h=\overline h\). Taking its bilinear scalar product with \(h\)
gives \(m_h=\|h\|^2\). Since \(p_h(v)<1\) everywhere, \(m_h<1\).
This proves \eqref{mf:stationarity}.

\subsection{An exact three-sector identity}

Suppress the index \(N\). For a retained set \(I\subset[N]\), put
\(J=I^c\), \(k=|I|\), and
\begin{align*}
 f_I(z)&=\prod_{i\in I}(v_i\cdot z),\\
 G_I(z)&=\sum_{\substack{i<j\\i,j\in I}}
        (v_i\wedge v_j)\prod_{b\in I\setminus\{i,j\}}(v_b\cdot z),\\
 H_{a,I}(z)&=\sum_{i\in I}\operatorname{sign}(a-i)\,v_i
                       \prod_{b\in I\setminus\{i\}}(v_b\cdot z),
                       \qquad a\in J.
\end{align*}
Empty products equal one and empty sums equal zero.

\begin{lemma}[All exterior sectors]\label{mf:sector-identity}
One has the exact identity
\begin{equation}\label{mf:three-sectors}
 T(A_N)=\sum_{I\subset[N]}\left(\frac tN\right)^{|I|}
 \left(
  \|G_I\|_{\mathrm F}^2+
  \sum_{a\in I^c}\|H_{a,I}\|_{\mathrm F}^2+
  \binom{|I^c|}{2}\|f_I\|_{\mathrm F}^2
 \right).
\end{equation}
\end{lemma}
\begin{proof}
Use Gram rows
\(w_i=\sqrt{t/N}\,v_i\oplus e_i\).
Split their wedge polynomial into the exterior sectors
\(\Lambda^2\mathbb C^r\), \(\mathbb C^r\wedge\mathbb C^N\), and
\(\Lambda^2\mathbb C^N\).
In the first sector, the coefficient of the added-variable monomial
\(\prod_{a\in J}x_a\) is \((t/N)^{k/2}G_I\).
In the second, with exterior index \(a\in J\), the coefficient of
\(\prod_{b\in J\setminus\{a\}}x_b\) is
\((t/N)^{k/2}H_{a,I}\), up to one irrelevant common sign.
In the third sector, each exterior pair in \(J\) gives
\((t/N)^{k/2}f_I\) multiplying the corresponding remaining monomial.
Distinct added-variable monomials are Fischer-orthogonal and each has
factorial weight one. The three exterior sectors are orthogonal.
Taking squared norms proves \eqref{mf:three-sectors}.
\end{proof}

For clarity, this also identifies precisely what diagonal dilution adds
to the derivative. Define
\[
 F(B)_{ij}=b_{ij}\per B(i\mid j),\qquad
 \operatorname{cut}_F(B[I];a)=
       \sum_{\substack{i<a<j\\i,j\in I}}F(B[I])_{ij},
\]
where retained indices keep their original labels in the cut.
The matrix \(F(B[I])\) has row sums \(\per B[I]\). Hence
\[
 \|H_{a,I}\|_{\mathrm F}^2
       =|I|\per B[I]-4\operatorname{cut}_F(B[I];a),
\]
and \eqref{mf:endpoint} gives
\begin{equation}\label{mf:derivative-dilution}
 D_1(A_N)=\sum_I\left(\frac tN\right)^{|I|}
 \left(D_1(B_N[I])
          +2\sum_{a\in I^c}\operatorname{cut}_F(B_N[I];a)\right).
\end{equation}
The cross-block cuts are nonnegative, as follows by differentiating
Lieb's cross-block scaling expansion, whose coefficients in the square
of the scaling variable are nonnegative \cite{Lieb1966}.

\subsection{Conditional Bernoulli moments}

For fixed \(z\), let \(I\) be a random retained set with independent
membership probabilities
\[
 p_i(z)=
 \frac{(t/N)|v_i\cdot z|^2}{1+(t/N)|v_i\cdot z|^2}.
\]
Summing \((t/N)^{|I|}|f_I(z)|^2\) produces exactly the product in
\eqref{mf:gaussian}. At \(z=\sqrt N\,h\), these probabilities are
\(p_h(v_i)\).

Away from the zero sets of the nonzero linear forms, let
\[
 U_i=\xi_i\,\frac{v_i}{\ell_h(v_i)},\qquad
 \xi_i\ \hbox{independent Bernoulli}(p_h(v_i)).
\]
Then
\begin{equation}\label{mf:bounded-moments}
 \mathbb E U_i=q_h(v_i),\qquad
 \mathbb E\|U_i\|^2
 =\frac{t\|v_i\|^2}{1+t|\ell_h(v_i)|^2}\le tL^2,\qquad
 \|\mathbb E U_i\|\le\frac{\sqrt t\,L}{2}.
\end{equation}
These quotient random variables are used only almost everywhere.
For each nonzero row its zero set is a complex hyperplane of Lebesgue
measure zero. Zero rows contribute zero and can be omitted from this
calculation. At a vanishing nonzero linear form, declaring \(U_i=0\)
would not reproduce the continuous limiting second moment in
\eqref{mf:bounded-moments}; that identification is not made.
Instead the moment expressions have continuous extensions in the original
polynomial integrals, and values on these null sets do not affect an
integral. The final functions \(p_h,q_h,\mathcal R\) have no quotient
singularities.

Set
\begin{align*}
 W_N(h)&=\frac1{N^2}\sum_{i<j}q_h(v_i)\wedge q_h(v_j),
 &Q_N(h)&=\frac1N\sum_iq_h(v_i),\\
 Q_{a,N}(h)&=\frac1N\sum_{i<a}q_h(v_i),
 &m_N(h)&=\frac1N\sum_ip_h(v_i).
\end{align*}
For the first exterior sector the random quotient is
\[
 \frac1N\sum_{i<j}U_i\wedge U_j,
\]
whose mean is \(NW_N\). In its variance, disjoint pairs have zero
covariance. There are \(O(N^3)\) pairs sharing an index and \(O(N^2)\)
identical pairs. Every surviving covariance is bounded by a constant
depending only on \(t,L\), using independence and
\eqref{mf:bounded-moments}. The factor \(N^{-2}\) thus makes the variance
\(O(N)\). The contribution of this sector divided by \(N^2\per A_N\)
equals
\begin{equation}\label{mf:first-sector-limit}
 \mathbb E_{\Pi_N}\|W_N(h)\|^2+O(N^{-1}).
\end{equation}

For a fixed missing index \(a\), the second-sector quotient is
\[
 \frac1{\sqrt N}\sum_{i\ne a}\operatorname{sign}(a-i)U_i.
\]
Its mean is
\(\sqrt N(2Q_{a,N}-Q_N+q_h(v_a)/N)\), and its variance is at most
\(tL^2\). Membership of \(a\) is independent of this sum. Summing over
\(a\) and using the uniform bounds gives the normalized second-sector
contribution
\begin{equation}\label{mf:second-sector-limit}
 \mathbb E_{\Pi_N}\frac1N\sum_a
 (1-p_h(v_a))\|Q_N(h)-2Q_{a,N}(h)\|^2+O(N^{-1}).
\end{equation}
The third sector counts missing pairs. Its normalized contribution is
\begin{equation}\label{mf:third-sector-limit}
 \frac12\mathbb E_{\Pi_N}(1-m_N(h))^2+O(N^{-1}).
\end{equation}
All three error estimates are uniform in \(h\).

\subsection{Convergence of the ordered functionals}

The discrete functions in
\eqref{mf:first-sector-limit}--\eqref{mf:third-sector-limit} converge to
their \(\mu\)-integral versions locally uniformly in \(h\).
For the order indicators, partition the \(x\)-interval into short
subintervals. The first marginal is Lebesgue measure, so boundaries
\(x=s\) and \(x=y\) have zero \(\mu\) and \(\mu\otimes\mu\) mass.
The uniform bound on \(q_h\) controls the contribution of the short
subintervals. This proves prefix convergence uniformly in \(s\).
A finite \(h\)-net and uniform continuity give joint local uniformity
in \(s,h\); the same argument handles the ordered double integral.

Thus \(\mathcal R\) is continuous and bounded, as are the discrete
integrands uniformly. Tightness of \(\Pi_N\), local uniform convergence,
and the three sector estimates give
\begin{equation}\label{mf:score-integral}
 \frac{T(A_N)}{\binom N2\per A_N}
       =\mathbb E_{\Pi_N}\mathcal R(h)+o(1).
\end{equation}
Concentration near \(K\) proves
\eqref{mf:criterion-bounds}. Equation~\eqref{mf:ratio} proves the
remaining assertions of Theorem~\ref{mf:criterion}.

\section{Proof of the subcritical obstruction}\label{mf:subcritical-proof}

\subsection{A dominated Gaussian limit}

Assume \(t\lambda_{\max}(\Sigma)<1\). Choose \(s>1\) and \(b<1\)
such that \(st\lambda_{\max}(\Sigma_N)\le b\) for all sufficiently
large \(N\). Since \(\log(1+x)\le x\),
\[
 \prod_i\left(1+\frac{st}{N}|v_i\cdot z|^2\right)
 \le e^{stz^*\Sigma_Nz}\le e^{b\|z\|^2}.
\]
This is integrable against the standard complex Gaussian, even after
multiplication by any fixed power of \(\|z\|\).
The bounded row assumption and the quadratic Taylor remainder show
pointwise convergence of the product with parameter \(t\) to
\(e^{tz^*\Sigma z}\). Dominated convergence gives
\[
 \per A_N\longrightarrow
 \mathbb E e^{tZ^*\Sigma Z}=\det(I-t\Sigma)^{-1},
\]
which proves \eqref{mf:permanent-limit}.

The positive principal-permanent weights
\[
 \mathbb P_N(I)=
 \frac{(t/N)^{|I|}\per B_N[I]}{\per A_N}
\]
define a random retained size \(K_N=|I|\). Its exponential moment is
\begin{equation}\label{mf:retained-mgf}
 \mathbb E s^{K_N}
 =\frac{\per(I+(st/N)B_N)}{\per A_N}\le C_s.
\end{equation}
Therefore every fixed moment of \(K_N\) is uniformly bounded.

\subsection{Polynomial bounds for compressed endpoint terms}

For any \(k\times k\) PSD Gram matrix \(B\), put \(C_k=\binom k2\).
Cauchy--Schwarz applied to the \(C_k\) wedge summands gives
\begin{align}
 T(B)
 &\le C_k\sum_{i<j}
       \det B[\{i,j\}]\per B[\{i,j\}^c]\notag\\
 &\le C_k^2\per B.\label{mf:wedge-bound}
\end{align}
The last inequality uses the elementary two-by-two inequality
\(\det\le\per\), followed by Lieb's principal-block inequality
\cite{Lieb1966}. The latter is also a consequence of the Fischer product
inequality for products of linear forms \cite{Reznick1993}.
It follows that
\begin{equation}\label{mf:derivative-bound}
 |D_1(B)|\le\frac{C_k+C_k^2}{2}\per B.
\end{equation}
The matrix \(F(B)\) is PSD: it is the Gram matrix of the vector-valued
polynomials \(v_i\prod_{j\ne i}\ell_j\).
Lieb's inequality gives
\(F(B)_{ii}=b_{ii}\per B(i\mid i)\le\per B\).
Thus \(|F(B)_{ij}|\le\per B\), and
\begin{equation}\label{mf:footrule-bound}
 \left|\sum_{i,j}|i-j|F(B)_{ij}\right|\le k^3\per B.
\end{equation}
The polynomial dependence on \(k\) is needed here; a factorial bound
would not suffice for \eqref{mf:retained-mgf}.

\subsection{Replacing inversions by the footrule}

Define
\begin{equation}\label{mf:footrule}
 \mathcal L(A)=\sum_{i,j}|i-j|a_{ij}\per A(i\mid j)
 =\sum_\sigma\left(\sum_i|i-\sigma(i)|\right)
                         \prod_i a_{i,\sigma(i)}.
\end{equation}
If a permutation fixes an index \(a\), deleting that index reduces its
inversion count by twice the number of edges crossing that cut in one
direction. It reduces its footrule by the same amount, because each
crossing edge has its length reduced by one and the crossing counts in
the two directions agree. Therefore
\[
 \operatorname{inv}(\sigma)-\sum_i|i-\sigma(i)|
\]
is unchanged by deleting fixed indices and compressing the remaining
order.
Expanding \(A_N=I+(t/N)B_N\) gives the exact identity
\begin{equation}\label{mf:compression}
 D_1(A_N)-\mathcal L(A_N)
 =\sum_I\left(\frac tN\right)^{|I|}
       \bigl(D_1(B_N[I])-\mathcal L(B_N[I])\bigr),
\end{equation}
where the submatrix quantities use their compressed order.
Equations~\eqref{mf:retained-mgf},
\eqref{mf:derivative-bound}, and \eqref{mf:footrule-bound} imply
\begin{equation}\label{mf:bounded-defect}
 |D_1(A_N)-\mathcal L(A_N)|\le C\per A_N
\end{equation}
with \(C\) independent of \(N\).

\subsection{The limiting cross trace}

For \(i\ne j\), Gaussian expansion of the cross permanent gives
\begin{equation}\label{mf:cross-minor}
 \per A_N(i\mid j)=\frac tN\,
 \mathbb E\left[
 (v_j\cdot Z)\overline{(v_i\cdot Z)}
 \prod_{a\ne i,j}\left(1+\frac tN|v_a\cdot Z|^2\right)\right].
\end{equation}
The domination used above, with the required Gaussian moment weight,
shows uniformly in \(i,j\) that the expectation in
\eqref{mf:cross-minor}, divided by \(\per A_N\), is
\(v_j^TR_0v_i+o(1)\).
Indeed, omitting two factors costs at most
\(O(N^{-1})\|Z\|^4\) times the dominating product; the full tilted
second-moment matrix converges in the fixed dimension.
Consequently
\[
 \frac{F(A_N)_{ij}}{\per A_N}
 =\frac{t^2}{N^2}(v_i^Tv_j)(v_j^TR_0v_i)+o(N^{-2})
\]
uniformly for \(i\ne j\). Weak convergence and \eqref{mf:footrule} yield
\begin{align}
 \frac{\mathcal L(A_N)}{N\per A_N}
 &\longrightarrow
 t^2\iint |x-y|(v^Tw)(w^TR_0v)\,d\mu(x,v)\,d\mu(y,w)\notag\\
 &=2t^2\int_0^1\operatorname{tr}
             \bigl(M_sR_0(\Sigma-M_s)\bigr)\,ds.\label{mf:footrule-limit}
\end{align}
The last equality integrates the two cut orientations of \(|x-y|\).
The two traces agree and are real, since \(R_0\) is a function of
\(\Sigma\).

\begin{lemma}[A resolvent cross-trace inequality]\label{mf:trace}
If \(A,B\succeq0\), \(S=A+B\), and \(t\|S\|<1\), then
\begin{equation}\label{mf:trace-inequality}
 \operatorname{tr}\bigl(A(I-tS)^{-1}B\bigr)
       \ge\operatorname{tr}(AB)\ge0.
\end{equation}
\end{lemma}
\begin{proof}
Restrict to the range of \(S\), since both \(A\) and \(B\) annihilate
its kernel. Write
\(A=S^{1/2}CS^{1/2}\), \(B=S^{1/2}(I-C)S^{1/2}\), with
\(0\preceq C\preceq I\), and diagonalize \(S\) with eigenvalues
\(\lambda_i>0\).
If \(f\ge0\) and \(x\mapsto xf(x)\) is nondecreasing on this spectrum,
direct expansion gives
\begin{align}
 \operatorname{tr}(Af(S)B)
 &=\sum_i\lambda_i^2f(\lambda_i)
                   \bigl(C_{ii}-(C^2)_{ii}\bigr)\notag\\
 &\quad+\frac12\sum_{i,j}(\lambda_i-\lambda_j)
       \bigl(\lambda_if(\lambda_i)-\lambda_jf(\lambda_j)\bigr)
                        |C_{ij}|^2\ge0.\label{mf:trace-expansion}
\end{align}
Every term is nonnegative, because \(C^2\preceq C\).
Apply this to \(f(x)=(1-tx)^{-1}-1\);
then \(xf(x)=tx^2/(1-tx)\) is nonnegative and increasing.
This proves the first inequality in \eqref{mf:trace-inequality};
the second is the standard trace inequality for two PSD matrices.
\end{proof}

The lower comparison is strictly positive after integration:
\begin{equation}\label{mf:strict-trace}
 2\int_0^1\operatorname{tr}\bigl(M_s(\Sigma-M_s)\bigr)\,ds
 =\iint|x-y||v^Tw|^2\,d\mu\,d\mu>0.
\end{equation}
The unweighted integral of \(|v^Tw|^2\) is
\(\operatorname{tr}(\Sigma^2)>0\), while \(x=y\) has zero
\(\mu\otimes\mu\) measure.
More quantitatively, \(|v^Tw|^2\le L^4\), and the mass of
\(|x-y|<\epsilon\) is at most \(2\epsilon\).
Taking \(\epsilon=\operatorname{tr}(\Sigma^2)/(4L^4)\) gives
\begin{equation}\label{mf:strict-lower-bound}
 \iint|x-y||v^Tw|^2\,d\mu\,d\mu
 \ge\frac{(\operatorname{tr}\Sigma^2)^2}{8L^4}.
\end{equation}
Equations~\eqref{mf:bounded-defect} and \eqref{mf:footrule-limit},
followed by Lemma~\ref{mf:trace}, prove
Theorem~\ref{mf:subcritical}.

\begin{proof}[Proof of Corollary~\ref{mf:arbitrary-orders}]
The ordered empirical measures lie in a compact space of probability
measures. Extract a weakly convergent subsequence from any proposed
violating subsequence. Its first marginal is Lebesgue measure and its
covariance is the corresponding limit of \(\Sigma_N\).
Apply Theorem~\ref{mf:subcritical} and
\eqref{mf:strict-lower-bound} on that subsequence.
This proves \eqref{mf:uniform-positive} and the version with uniform
\(\varepsilon,s\).
\end{proof}

\section{The exact spherical saddles and a stronger Gini bound}\label{mf:sphere}

\subsection{The anisotropic spherical potential}

Let \(Y\) be uniform on the real sphere \(S^2\) and
\[
 v=T_\tau Y,\qquad T_\tau=\operatorname{diag}(1,1,\tau),
 \qquad 0<\tau<1.
\]
The marginal law of \(v\) has covariance
\(\Sigma=\operatorname{diag}(1/3,1/3,\tau^2/3)\).
The order is specified below and does not affect \(\Phi\).

\begin{proposition}[Exact coherent transition]\label{mf:spherical-saddles}
For this spherical law, \(h=0\) is the unique maximizer of \(\Phi\)
when \(t\le3\). When \(t>3\), its maximizers are exactly
\begin{equation}\label{mf:circular-poles}
 h=e^{i\phi}\sqrt{m/2}(1,i,0)
 \quad\hbox{or}\quad
 h=e^{i\phi}\sqrt{m/2}(1,-i,0),
\end{equation}
where \(m\in(0,1)\) is uniquely determined by
\begin{equation}\label{mf:radial-saddle}
 m=m(k):=\int_0^1\frac{k(1-z^2)}{1+k(1-z^2)}\,dz,\qquad
 k=\frac{tm}{2}.
\end{equation}
Equivalently, \(t=2k/m(k)\), a strictly increasing function of \(k>0\),
and
\begin{equation}\label{mf:mass-formula}
 1-m(k)=
 \frac{\operatorname{artanh}\sqrt{k/(1+k)}}{\sqrt{k(1+k)}}.
\end{equation}
\end{proposition}
\begin{proof}
Multiply \(h\) by a phase so that the real and imaginary parts of
\(T_\tau h\) are orthogonal, and rotate \(Y\).
If their squared lengths are \(\lambda_1,\lambda_2\), the logarithmic
term in \(\Phi\) becomes
\[
 \mathbb E\log\bigl(1+t(\lambda_1Y_1^2+\lambda_2Y_2^2)\bigr).
\]
It is symmetric and strictly concave along each line with fixed positive
\(\lambda_1+\lambda_2\): the second derivative is the negative
expectation of a nonzero square proportional to
\((Y_1^2-Y_2^2)^2\).
It is therefore maximized at \(\lambda_1=\lambda_2\), and is strictly
increasing in their sum. That sum is
\(\|T_\tau h\|^2\le\|h\|^2\), with equality only if \(h_3=0\).
Thus, at fixed \(\|h\|^2=m>0\), the maximizing directions are precisely
\eqref{mf:circular-poles}.

The remaining radial function is strictly concave:
\[
 \varphi_t(m)=\int_0^1\log(1+(tm/2)(1-z^2))\,dz-m.
\]
Its derivative at zero is \(t/3-1\). It tends to \(-\infty\) as
\(m\to\infty\), so the asserted unique radial maximum follows.
The stationary equation gives \eqref{mf:radial-saddle}; its right side
is strictly between zero and one.
The function \(m(k)/k\) is strictly decreasing, proving the asserted
monotonicity of \(t(k)\). Integrating
 \((1+k-kz^2)^{-1}\) gives \eqref{mf:mass-formula}.
\end{proof}

The transition at \(t=3\) is a transition of the coherent maximizers.
It does not by itself imply a negative endpoint derivative for every
\(t>3\); the ordered score must still exceed one.

\subsection{A quantitative strict Gini inequality}

Let \(R\) have density \(1/[2(1+r^2)^{3/2}]\), independently of a
uniform angle \(\Theta\). For an integrable real random variable \(Z\),
write \(\mathcal C(Z)=\tfrac12\mathbb E|Z-Z'|\), with an independent
copy \(Z'\). Put \(X=\sqrt2R\cos\Theta\).

\begin{lemma}[A stronger spherical Gini gain]\label{mf:strong-gini}
One has
\begin{equation}\label{mf:gini-gain}
 \mathcal C(X)\ge\frac1{\sqrt2}
             +\frac{\pi^2-8}{32\sqrt2\,\pi}
             >\frac1{\sqrt2}+\frac{13}{1000}.
\end{equation}
\end{lemma}
\begin{proof}
Let \(U=|R|\) and
\(g(a,b)=\mathbb E\max(aU,bU')\) for \(a,b\ge0\).
Independent symmetric signs show that
\(g(a,b)=\mathbb E|aR-bR'|\).
The distribution function of \(U\) is \(u/\sqrt{1+u^2}\), so
\[
 g(a,b)=\int_0^\infty
 \left(1-\frac{s^2}{\sqrt{(s^2+a^2)(s^2+b^2)}}\right)\,ds.
\]
For positive \(a,b\), differentiating the expectation of the maximum
gives
\[
 g_{aa}(a,b)=\frac1b\int_0^\infty
 \frac{u^2\,du}
 {(1+u^2)^{3/2}(1+(a/b)^2u^2)^{3/2}}.
\]
Differentiation is justified locally by an integrable majorant when
\(a,b\) range in a compact subset of \((0,\infty)\).
Homogeneity gives \(g_{ab}=-(a/b)g_{aa}\) and
\(g_{bb}=(a/b)^2g_{aa}\).

Set \(H(x)=g(1+x,1-x)\) and \(a=1+x,b=1-x\).
Substitution \(u=s\sqrt{b/a}\) yields
\[
 H''(x)=\frac4{(ab)^{3/2}}\int_0^\infty
 \frac{s^2\,ds}
 {\bigl((1+s^2)^2+4(1-ab)s^2/(ab)\bigr)^{3/2}}.
\]
Because \(4s^2\le(1+s^2)^2\) and \(0<ab\le1\), the denominator
inside the parentheses is at most \((1+s^2)^2/(ab)\).
Hence
\[
 H''(x)\ge4\int_0^\infty\frac{s^2}{(1+s^2)^3}\,ds
          =\frac\pi4.
\]
Symmetry gives \(H'(0)=0\).
The radial distribution function
\(F_R(r)=\tfrac12+r/(2\sqrt{1+r^2})\) gives
\[
 \mathcal C(R)=\int_{\mathbb R}F_R(r)(1-F_R(r))\,dr
             =\int_{\mathbb R}\frac{dr}{4(1+r^2)}=\frac\pi4,
\]
so \(H(0)=\pi/2\).
Integrating twice and rescaling proves
\begin{equation}\label{mf:uniform-convexity}
 g(a,b)\ge\frac\pi4(a+b)
               +\frac\pi{16}\frac{(a-b)^2}{a+b}.
\end{equation}
Continuity extends this to the boundary, interpreting the final term as
zero at \(a=b=0\). For example, continuity follows from
\(\mathbb E U=1\) and the Lipschitz bound for the maximum.

By symmetry \(X\) has the law of \(SR\), with
\(S=\sqrt2|\cos\Theta|\) independent of \(R\).
Here
\[
 \mathbb ES=\frac{2\sqrt2}{\pi},\qquad
 \mathbb ES^2=1,\qquad 0\le S\le\sqrt2.
\]
Average \eqref{mf:uniform-convexity} over \(S,S'\). Since
\(S+S'\le2\sqrt2\) and
\(\mathbb E(S-S')^2=2\operatorname{Var}(S)\), this gives
\[
 \mathcal C(X)=\tfrac12\mathbb Eg(S,S')
 \ge\frac\pi4\mathbb ES+\frac{\pi}{32\sqrt2}\operatorname{Var}(S)
 =\frac1{\sqrt2}+\frac{\pi^2-8}{32\sqrt2\,\pi}.
\]
For the displayed rational bound, use
\(\pi>157/50\), \(\pi<22/7\), and \(\sqrt2<99/70\).
The resulting rational lower bound for the gain exceeds \(13/1000\).
\end{proof}

\section{A bounded-condition consequence}\label{mf:concrete}

Use spherical coordinates
\[
 Y=(\sqrt{1-z^2}\cos\theta,\sqrt{1-z^2}\sin\theta,z),
\]
where \(z\) is uniform on \([-1,1]\) and \(\theta\) is independently
uniform on \([0,2\pi]\). Order rows by the real quotient
\begin{equation}\label{mf:sort-coordinate}
 X_\tau(Y)=\frac{\sqrt2\,\tau z\cos\theta}{\sqrt{1-z^2}}.
\end{equation}
Its values at the two sphere poles can be assigned arbitrarily, since
those points have measure zero.

At the positive saddle in \eqref{mf:circular-poles}, write \(a=\sqrt m\)
and use the unitary coordinates with first basis vector
\((1,i,0)/\sqrt2\). Dividing a row form by its pole value gives
\[
 c=(1,e^{-2i\theta},
          \sqrt2\,\tau z e^{-i\theta}/\sqrt{1-z^2}).
\]
In these coordinates \(q_h=(p/a)c\), with
\[
 p(z)=\frac{k(1-z^2)}{1+k(1-z^2)}.
\]
Consequently the real part of one exterior component is
\begin{equation}\label{mf:effective-gini}
 C_{\mathrm{eff}}:=\Re(W_{02})
 =\frac1m\cdot\frac12
      \mathbb E\bigl[p(z)p(z')|X_\tau-X'_\tau|\bigr].
\end{equation}
Although the quotient is unbounded, \(p\,c\) is bounded. Equivalently,
\(q_h\) is a continuous bounded function on the original compact sphere.
No logarithmic clipping or radial truncation is needed.

\begin{lemma}[Loss from diagonal dilution]\label{mf:gini-loss}
With \(C_X=\mathcal C(\sqrt2R\cos\Theta)\) as above,
\begin{equation}\label{mf:gini-loss-bound}
 C_{\mathrm{eff}}\ge
 \tau\left(C_X-\frac{\sqrt2}{\sqrt k}
                   -\frac{\log(1+k)}k\right).
\end{equation}
\end{lemma}
\begin{proof}
Let \(C_p=\tfrac12\mathbb E[pp'|X_\tau-X'_\tau|]\).
The inequality \(1-pp'\le(1-p)+(1-p')\), followed by the triangle
inequality, gives
\[
 \tau C_X-C_p\le
 \mathbb E[(1-p)|X_\tau|]+(1-m)\mathbb E|X_\tau|.
\]
Direct integration gives
\[
 \mathbb E|X_\tau|=\frac{2\sqrt2\,\tau}{\pi},\qquad
 \mathbb E[(1-p)|X_\tau|]
 =\frac{2\sqrt2\,\tau}{\pi}
                  \frac{\arctan\sqrt k}{\sqrt k}.
\]
For the second identity substitute \(u=\sqrt{1-z^2}\) into the
radial integral. Also
\[
 1-m=\int_0^1\frac{dz}{1+k(1-z^2)}
 \le\int_0^1\frac{dz}{1+k(1-z)}
 =\frac{\log(1+k)}k.
\]
Now use \(\arctan\sqrt k\le\pi/2\),
\(2\sqrt2/\pi<1\), and \(C_{\mathrm{eff}}=C_p/m\ge C_p\).
\end{proof}

\begin{theorem}[An eventual family with condition number below \(20002\)]
\label{mf:20002}
There is an explicit algebraic sequence of real positive-definite
correlation matrices \(C_N\), defined for every sufficiently large order,
such that
\begin{equation}\label{mf:strong-derivative}
 \limsup_{N\to\infty}
 \frac{D_1(C_N)}{N(N-1)\per C_N}<-\frac1{800}.
\end{equation}
Its off-diagonal entries have modulus at most \(60000/N\), its least
eigenvalue tends to one, and its condition number tends to \(20001\).
In particular its condition number is eventually less than \(20002\).
No explicit sufficient order is asserted in this theorem.
\end{theorem}
\begin{proof}
Take
\[
 \tau=\frac{999}{1000},\qquad t=60000.
\]
The saddle parameter satisfies \(k>25000\).
Indeed \(t(k)=2k/m(k)\) is strictly increasing, and
\[
 1-m(25000)\le\frac{\log25001}{25000}<\frac16
\]
implies \(t(25000)<60000\).
Both loss functions in \eqref{mf:gini-loss-bound} decrease with \(k\).
At \(k=25000\), the bounds
\(\sqrt2<3/2\), \(\sqrt{25000}>158\), and \(\log25001<11\) give
\[
 \frac{\sqrt2}{\sqrt k}+\frac{\log(1+k)}k
 <\frac3{316}+\frac{11}{25000}<\frac1{100}.
\]
For example, \(\log25001<11\) follows from \(e>8/3\) and
\((8/3)^{11}>25001\).
Lemmas~\ref{mf:strong-gini} and \ref{mf:gini-loss} therefore imply
\[
 C_{\mathrm{eff}}>
 \tau\left(\frac1{\sqrt2}+\frac3{1000}\right)
 >\frac1{\sqrt2}+\frac2{1000}.
\]
Using \(\sqrt2>5/4\) shows
\[
 2C_{\mathrm{eff}}^2>1+\frac1{200}.
\]
The two saddle circles are conjugate, so their full score is constant.
Theorem~\ref{mf:criterion}, even using only the first exterior sector,
proves \eqref{mf:strong-derivative} for any ordered discretization
converging to this sphere law.

Here is one deterministic algebraic discretization in every order.
Let \(m=\lfloor\sqrt N\rfloor\), and form the \(m^2\) sphere rows
\[
 z_\ell=\frac{2\ell-1}{m}-1,\quad
 \theta_j=\frac{2\pi j}{m},\quad
 Y_{j,\ell}=
  (\sqrt{1-z_\ell^2}\cos\theta_j,
   \sqrt{1-z_\ell^2}\sin\theta_j,z_\ell),
\]
for \(1\le\ell\le m\), \(0\le j<m\).
Apply \(T_\tau\), add \(N-m^2\) copies of \((1,0,0)\), and sort by
\eqref{mf:sort-coordinate}, breaking ties lexicographically by the
original labels. The grid avoids the two poles and all entries are
algebraic.

These empirical sphere measures converge to the uniform sphere law.
The quotient is continuous away from the two poles and its limiting
distribution has no atoms. Its empirical distribution functions therefore
converge uniformly to its continuous distribution function \(F_X\), and
the largest relative tie multiplicity tends to zero.
The sorted empirical position differs uniformly by \(o(1)\) from
\(F_X(X_\tau(Y))\). Hence the joint ordered law converges to that of
\((F_X(X_\tau(Y)),T_\tau Y)\). Its first marginal is uniform on
\([0,1]\), verifying \eqref{mf:empirical}.

Set \(A_N=I+(60000/N)V_NV_N^T\) and normalize its diagonal:
\[
 C_N=\Delta_N^{-1/2}A_N\Delta_N^{-1/2},\qquad
 \Delta_N=\operatorname{diag}(A_N).
\]
Every permutation monomial, and hence the permanent and its derivative,
is multiplied by the same positive factor. Thus the normalized derivative
ratio is unchanged.
The covariance eigenvalues converge to \(1/3,1/3,\tau^2/3\), so the
condition number of \(A_N\) tends to \(1+60000/3=20001\).
Its least eigenvalue is one for large \(N\).
Since \(\Delta_N=I+O(N^{-1})\) in operator norm and \(\|A_N\|\) is
bounded, \(\|C_N-A_N\|\to0\). This proves the asserted eigenvalue and
condition-number limits. The unit bound on \(\|T_\tau Y\|\), and
diagonal normalization factors at least one, give the off-diagonal bound.
\end{proof}

Before correlation normalization, exactly \(N-3\) eigenvalues equal one
when the rows span \(\mathbb R^3\). Nonconstant diagonal normalization
need not preserve that exact multiplicity; only the stated limiting
spectral conclusions are claimed for \(C_N\).
This theorem concerns large orders only and supplies no explicit
sufficient order; the finite certificate of Section~\ref{sec:certificate}
is a separate construction.

\section{A four-dimensional sphere and condition numbers below 257}
\label{mf:four-sphere}

The second exterior sector can make a decisive contribution. We now use
the complete score \eqref{mf:score} for a four-dimensional real spherical
law. The resulting condition-number bound is much smaller, while the
preceding construction retains its advantage of using three-dimensional
latent rows.

Let \(Y\) be uniform on \(S^3\subset\mathbb R^4\), and set
\[
 v=T_\tau Y,\qquad T_\tau=\operatorname{diag}(1,1,\tau,\tau),
 \qquad 0<\tau<1.
\]
Write
\[
 w_1=Y_1+iY_2,\qquad w_2=Y_3+iY_4,\qquad
 S=|w_1|^2,\qquad b=w_2/w_1.
\]
The variable \(S\) is uniform on \([0,1]\), and the two complex phases
are independent and uniform. These statements follow, for example, by
normalizing two independent standard complex Gaussians.
The ratio \(b\) has density
\begin{equation}\label{mf:complex-cauchy}
 \frac{1}{\pi(1+|b|^2)^2}\,d^2b.
\end{equation}
Indeed \(S=(1+|b|^2)^{-1}\), and changing from the uniform \(S\) to the
radius of \(b\) gives the radial density \(2r/(1+r^2)^2\).
Order the rows by \(x=\Re b\); the exceptional set \(w_1=0\) has
measure zero.

\subsection{Exact saddle geometry and marginal formulas}

The proof of Proposition~\ref{mf:spherical-saddles} applies in real
dimension four: rotational invariance and strict concavity in the two
squared singular values force a maximizing complex direction to be
circular in the first coordinate plane.
For \(t>4\), the maximizers are exactly
\[
 h=e^{i\phi}\sqrt{m/2}(1,i,0,0)
 \quad\hbox{or its conjugate},
\]
where
\begin{equation}\label{mf:four-saddle}
 m=\int_0^1\frac{kS}{1+kS}\,dS
   =1-\frac{\log(1+k)}k,\qquad k=\frac{tm}{2}.
\end{equation}
For completeness, the radial potential is
\(\int_0^1\log(1+(tm/2)S)\,dS-m\); it is strictly concave, has
derivative \(t/4-1\) at zero, and tends to \(-\infty\).
This proves uniqueness of the positive radial saddle.
The function \(t(k)=2k/m(k)\) is strictly increasing by the same
integral argument as before.

Put \(K=1+k\), \(\delta=1-m=\log K/k\), and
\[
 p=\frac{kS}{1+kS}=\frac{k}{K+|b|^2}.
\]
The unweighted marginal density of \(x=\Re b\), and its density
weighted by \(p\), are respectively
\begin{align}
 f(x)&=\frac1{2(1+x^2)^{3/2}},\label{mf:four-f}\\
 f_p(x)&=
 \frac1{2(1+x^2)^{3/2}}
 -\frac1{k\sqrt{1+x^2}}
 +\frac1{k\sqrt{K+x^2}}.\label{mf:four-fp}
\end{align}
To check the second formula, use
\[
 \frac{k}{u^2(u+k)}=\frac1{u^2}-\frac1{ku}
                         +\frac1{k(u+k)}
\]
with \(u=1+|b|^2\), and integrate the imaginary coordinate.
Thus \(f_p\ge0\), \(f-f_p\ge0\), and their masses are \(m\) and
\(\delta\), respectively.

Define the centered prefix mass and first moment
\[
 \mathcal A(x)=\int_{-\infty}^x f_p(u)\,du-\frac m2,\qquad
 \mathcal M(x)=\int_{-\infty}^x u f_p(u)\,du.
\]
Elementary antiderivatives give
\begin{align}
 \mathcal A(x)
 &=\frac{x}{2\sqrt{1+x^2}}
   -\frac{\operatorname{arsinh}x-
          \operatorname{arsinh}(x/\sqrt K)}k,\label{mf:four-prefix}\\
 \mathcal M(x)
 &=-\frac1{2\sqrt{1+x^2}}
   +\frac1{\sqrt{K+x^2}+\sqrt{1+x^2}}.\label{mf:four-moment}
\end{align}
For the second identity, the difference
\((\sqrt{K+x^2}-\sqrt{1+x^2})/k\) is rationalized.
The integration constant is fixed by the zero limit at \(-\infty\).

\subsection{The full score in one-dimensional integrals}

Use the unitary basis
\((1,\pm i,0,0)/\sqrt2\), \((0,0,1,\pm i)/\sqrt2\).
If \(\theta=\arg w_1\) and \(a=e^{-2i\theta}\), the row divided by
its first pole value has coefficients
\begin{equation}\label{mf:four-coefficients}
 c=(1,a,\tau b,\tau a\overline b).
\end{equation}
The phase \(a\) is uniform and independent of \(b\). Conditional on
\(x=\Re b\), the imaginary part of \(b\) has a symmetric law, also
after multiplication by \(p\).
Consequently only coordinates zero and two of the averaged vector
\(q_h=(p/\sqrt m)c\) survive, with densities
\[
 \frac1{\sqrt m}\bigl(f_p(x),0,\tau x f_p(x),0\bigr).
\]
Writing
\begin{equation}\label{mf:four-gamma}
 \Gamma_k=\frac1{2m}\iint|x-y|f_p(x)f_p(y)\,dx\,dy,
\end{equation}
we obtain \(\|W_h\|^2=\tau^2\Gamma_k^2\).
The prefix vector gives
\[
 \|Q_h-2Q_h(x)\|^2
   =\frac4m\bigl(\mathcal A(x)^2+\tau^2\mathcal M(x)^2\bigr).
\]
It follows that the exact limiting score at either saddle circle is
\begin{equation}\label{mf:four-full-score}
 \mathcal R_*=
 2\tau^2\Gamma_k^2+
 \frac8m\int_{\mathbb R}(f-f_p)
         \bigl(\mathcal A^2+\tau^2\mathcal M^2\bigr)\,dx+
 \delta^2.
\end{equation}
All terms are finite. This calculation uses the complete
Theorem~\ref{mf:criterion}; omitting the middle term would not give the
bound below.

\begin{lemma}[Bounds for the two nonconstant sectors]\label{mf:four-bounds}
One has
\begin{equation}\label{mf:four-gamma-lower}
 \Gamma_k\ge
 \frac{\pi/4-2\sqrt K/k}{m}.
\end{equation}
If
\[
 \rho_k=\frac12-\frac\delta2-\frac1{\sqrt K+1}>0,
\]
then
\begin{equation}\label{mf:four-mixed-lower}
 \frac8m\int(f-f_p)(\mathcal A^2+\tau^2\mathcal M^2)\,dx
 \ge\frac{8\tau^2\delta}{m}\rho_k^2.
\end{equation}
\end{lemma}
\begin{proof}
For the probability density \(f\), the absolute-difference potential is
\[
 g(x)=\int|x-y|f(y)\,dy=\sqrt{1+x^2}.
\]
Its derivative is \(2F(x)-1=x/\sqrt{1+x^2}\), and its value at zero
is one, proving the formula. The half mean absolute difference of this
law is \(\pi/4\), by the radial calculation in
Lemma~\ref{mf:strong-gini}.
If \(d=f-f_p\), expansion of \(f_p=f-d\) yields
\[
 \frac12\iint|x-y|f_p(x)f_p(y)\,dx\,dy
 \ge\frac\pi4-\int d(x)g(x)\,dx.
\]
The omitted term is
\(\tfrac12\iint|x-y|d(x)d(y)\,dx\,dy\ge0\).
Using \eqref{mf:four-fp},
\begin{align*}
 \int d(x)g(x)\,dx
 &=\frac1k\int_{\mathbb R}
       \left(1-\frac{\sqrt{1+x^2}}{\sqrt{K+x^2}}\right)\,dx\\
 &\le\frac2k\int_0^\infty
       \left(1-\frac{x}{\sqrt{K+x^2}}\right)\,dx
 =\frac{2\sqrt K}{k}.
\end{align*}
Dividing by \(m\) proves \eqref{mf:four-gamma-lower}.

For the mixed term, compare with the unweighted pair
\[
 \mathcal A_\infty(x)=\frac{x}{2\sqrt{1+x^2}},\qquad
 \mathcal M_\infty(x)=-\frac1{2\sqrt{1+x^2}},
\]
which satisfies \(\mathcal A_\infty^2+\mathcal M_\infty^2=1/4\).
Equations~\eqref{mf:four-prefix}--\eqref{mf:four-moment} give the
uniform estimates
\[
 |\mathcal A-\mathcal A_\infty|\le\delta/2,\qquad
 |\mathcal M-\mathcal M_\infty|\le\frac1{\sqrt K+1}.
\]
The first follows also by integrating the even nonnegative density
\(f-f_p\) of mass \(\delta\) from zero to \(x\).
The triangle inequality in \(\mathbb R^2\) therefore gives
\[
 \mathcal A^2+\tau^2\mathcal M^2
 \ge\tau^2(\mathcal A^2+\mathcal M^2)
 \ge\tau^2\rho_k^2.
\]
Integrating against \(f-f_p\), whose mass is \(\delta\), proves
\eqref{mf:four-mixed-lower}.
\end{proof}

\begin{theorem}[An eventual family with condition number below \(257\)]
\label{mf:257}
There is an explicit algebraic sequence of real positive-definite
correlation matrices \(E_N\), defined for every sufficiently large order,
such that
\begin{equation}\label{mf:four-derivative}
 \lim_{N\to\infty}
 \frac{D_1(E_N)}{N(N-1)\per E_N}<-\frac1{1000}.
\end{equation}
Its off-diagonal entries have modulus at most \(1020/N\), its least
eigenvalue tends to one, and its condition number tends to \(256\).
In particular its condition number is eventually below \(257\).
The unnormalized mean-field perturbation has latent rank four.
This theorem asserts no explicit sufficient order.
\end{theorem}
\begin{proof}
Set \(\tau=999/1000\) and \(t=1020\).
Equation~\eqref{mf:four-saddle} implies \(k<510\).
Since \(m(500)>1-7/500>50/51\), monotonicity of \(t(k)\) gives
\(t(500)<1020\), and therefore \(k>500\).
On this interval,
\[
 6<\log(1+k)<7,\qquad
 \frac6{510}<\delta<\frac7{500},\qquad
 22<\sqrt K<\frac{1131}{50}.
\]
For example, the logarithmic bounds follow from
\(e<11/4\), \((11/4)^6<501\), \(e>8/3\), and
\((8/3)^7>511\). These are strict elementary inequalities.
Using \(\pi>157/50\), Lemma~\ref{mf:four-bounds} gives
\begin{align*}
 \Gamma_k
 &>\frac{157/200-2(1131/50)/500}{1-6/510}
   =\frac{295171}{420000},\\
 \rho_k
 &>\frac12-\frac7{1000}-\frac1{23}
   =\frac{10339}{23000}>0.
\end{align*}
In \eqref{mf:four-full-score}, use \(1/m>1\) and
\(\delta^2\ge\tau^2\delta^2\). The resulting entirely rational bound is
\begin{align}
 \mathcal R_*
 &>\left(\frac{999}{1000}\right)^2
 \left[
 2\left(\frac{295171}{420000}\right)^2
 +8\frac6{510}\left(\frac{10339}{23000}\right)^2
 +\left(\frac6{510}\right)^2
 \right]\notag\\
 &=\frac{1505674035849200604129}
         {1498233800000000000000}
 >1+\frac1{250}.\label{mf:four-rational-margin}
\end{align}
The full score is constant on the two conjugate saddle circles.
Theorem~\ref{mf:criterion} and \eqref{mf:ratio} prove
\eqref{mf:four-derivative}.

An explicit discretization uses \(m=\lfloor N^{1/3}\rfloor\) and the
\(m^3\) rows
\[
 S_\ell=\frac{2\ell-1}{2m},\quad
 \theta_j=\frac{2\pi j}{m},\quad \psi_a=\frac{2\pi a}{m},
\]
\[
 Y_{j,a,\ell}=
 (\sqrt{S_\ell}\cos\theta_j,\sqrt{S_\ell}\sin\theta_j,
  \sqrt{1-S_\ell}\cos\psi_a,\sqrt{1-S_\ell}\sin\psi_a).
\]
Here \(1\le\ell\le m\) and \(0\le j,a<m\).
Apply \(T_\tau\), pad to \(N\) rows with \((1,0,0,0)\), and sort by
\[
 \Re b=\sqrt{\frac{1-S_\ell}{S_\ell}}\cos(\psi_a-\theta_j)
\]
with a fixed lexicographic tie convention. All entries are algebraic
and \(w_1\ne0\) on the grid. The number of added rows is \(O(N^{2/3})\).
The unweighted empirical sphere measures converge to the uniform law.
The limiting ordering coordinate has the continuous density \(f\) in
\eqref{mf:four-f}; the same distribution-function argument used in
Theorem~\ref{mf:20002} verifies ordered empirical convergence.

Set \(A_N=I+(1020/N)V_NV_N^T\) and normalize its diagonal as before.
Its covariance tends to
\(\operatorname{diag}(1/4,1/4,\tau^2/4,\tau^2/4)\), so its condition
number tends to \(1+1020/4=256\).
The row norms are at most one. Diagonal normalization preserves the
derivative ratio, changes the operator by \(o(1)\), and gives the claimed
off-diagonal bound and spectral limits.
\end{proof}

Theorems~\ref{mf:20002} and \ref{mf:257} express a tradeoff: the first
uses three-dimensional latent vectors and has the larger derivative
margin, while the second uses four-dimensional latent vectors and has
the much smaller condition-number bound.
Both are eventual-order conclusions from conventional asymptotic proofs.
Neither claims a finite negative Lean certificate.


% Input immediately after mean_field.tex, before spectral_obstruction.tex.
% A scalar refinement; the original mf:257 family remains unchanged.

\subsection{A refinement to limiting condition number 227}
\label{rf:subsection}

A different choice of the two scalar parameters improves the conditioning
bound using the same spherical law, ordering, and sector estimates.

\begin{theorem}[An eventual family with condition number below \(228\)]
\label{rf:228}
There is an explicit algebraic sequence of real positive-definite
correlation matrices \(\widetilde E_N\), defined in every sufficiently
large order, such that
\[
 \lim_{N\to\infty}
 \frac{D_1(\widetilde E_N)}{N(N-1)\per\widetilde E_N}
 <-\frac1{40000},\qquad
 |(\widetilde E_N)_{ij}|\le\frac{904}{N}\quad(i\ne j).
\]
Its least eigenvalue tends to one, and its condition number tends to
\(227\), so it is eventually below \(228\). Before correlation
normalization it is a mean-field perturbation with latent rank four.
No explicit sufficient order is asserted.
\end{theorem}
\begin{proof}
Use the construction of Theorem~\ref{mf:257}, replacing its parameters
by \(\tau=9999/10000\) and \(t=904\). Write
\(m(k)=1-\log(1+k)/k\) and \(t(k)=2k/m(k)\).
The elementary bounds \(5/2<e<11/4\),
\((11/4)^6<445\), and \((5/2)^7>447\) give
\(6<\log445<\log447<7\). Consequently
\[
 t(444)<\frac{2\cdot444^2}{437}
       =\frac{394272}{437}<904,
 \qquad
 t(446)>\frac{2\cdot446^2}{440}
       =\frac{49729}{55}>904.
\]
Strict monotonicity of \(t(k)\) proves \(444<k<446\).
With \(K=1+k\) and \(\delta=1-m\), this implies
\[
 \frac6{446}<\delta<\frac7{444},\qquad
 \frac{2109}{100}<\sqrt K<\frac{423}{20}.
\]
The square-root inequalities follow by squaring:
\((2109/100)^2<445\) and \((423/20)^2>447\).
Using \(\pi>31415/10000\) in Lemma~\ref{mf:four-bounds} gives
\begin{align*}
 \Gamma_k
 &>\frac{31415/40000-2(423/20)/444}{1-6/446}
   =\frac{45552433}{65120000}=:g>0,\\
 \rho_k
 &>\frac12-\frac7{888}-\frac{100}{2209}
   =\frac{876533}{1961592}=:r>0.
\end{align*}
As in \eqref{mf:four-rational-margin}, use \(1/m>1\) and
\(\delta^2\ge\tau^2\delta^2\) in the full score. Exact rational
arithmetic yields
\begin{align}
 \mathcal R_*
 &>\left(\frac{9999}{10000}\right)^2
  \left[2g^2+8\frac3{223}r^2+
                     \left(\frac3{223}\right)^2\right]\notag\\
 &=\frac{425269160832679340834525001397641}
         {425220873828711680000000000000000}
 >1+\frac1{10000}.\label{rf:rational-margin}
\end{align}
Theorem~\ref{mf:criterion} and \eqref{mf:ratio} give the stated
derivative limit.

Use the same algebraic sphere grid, padding, and sorting as in
Theorem~\ref{mf:257}, with the new \(\tau\), and set
\(\widetilde A_N=I_N+(904/N)V_NV_N^T\). The row norms are at most
one, and the limiting covariance is
\(\operatorname{diag}(1/4,1/4,\tau^2/4,\tau^2/4)\).
The condition number therefore tends to \(1+904/4=227\).
Correlation normalization preserves the normalized derivative, changes
the operator by \(o(1)\), and gives the off-diagonal bound \(904/N\)
and the least-eigenvalue limit as before.
\end{proof}

The family in Theorem~\ref{mf:257} retains the larger proved derivative
margin \(1/1000\). The present theorem improves its condition-number
bound with a smaller proved margin; it is a scalar refinement of the
same analytic construction.

% Input-ready conventional-proof component.  Labels use the prefix sp:.
% Requires the mean_field.tex component and its mf: labels.

\section{A spectral obstruction with an explicit positive margin}
\label{sp:section}

The full score in Theorem~\ref{mf:criterion} admits an operator
factorization that gives an obstruction in every fixed latent dimension.
Throughout this section the rows and empirical measures satisfy the
bounded-row hypotheses of Section~\ref{mf:setup}. Eigenvalues of the
limiting covariance \(\Sigma\) are listed in decreasing order, with
missing eigenvalues interpreted as zero. All analytic arguments in
this section are conventional proofs.

\begin{theorem}[At most two supercritical covariance directions]
\label{sp:main}
Let \(r\ge1\), \(t>0\), and \(\|v_{i,N}\|\le L\), and suppose
\[
 \Sigma_N=\frac1N V_N^TV_N\longrightarrow\Sigma,\qquad
 t\lambda_1(\Sigma)>1,\qquad t\lambda_3(\Sigma)\le1.
\]
For \(A_N=I_N+(t/N)V_NV_N^T\), every sequence of row orders satisfies
\begin{equation}\label{sp:positive-margin}
 \liminf_{N\to\infty}
 \frac{D_1(A_N)}{N(N-1)\per A_N}
 \ge
 \frac{\bigl(t\lambda_1(\Sigma)-1\bigr)^2}{64r\,t^2L^4}>0.
\end{equation}
In particular the endpoint derivative is eventually positive.
The hypothesis includes the boundary \(t\lambda_3(\Sigma)=1\).
\end{theorem}

The nontrivial eigenvalues of \(A_N\) converge to
\(1+t\lambda_j(\Sigma)\). Thus the theorem permits at most two
limiting eigenvalues strictly above two, while requiring the largest
one to exceed two. This is a statement about the specified bounded-row,
fixed-rank mean-field families. Its conclusion does not assert a
spectral criterion for unrestricted finite matrices.

\subsection{The full score as a Hilbert--Schmidt norm}

Fix an ordered limit \(\mu\) as in \eqref{mf:empirical} and a point
\(h\in\mathbb C^r\). Use \(p=p_h\), \(q=q_h\), and \(m=m_h\)
from \eqref{mf:p}--\eqref{mf:q}, and define the real symmetric matrix
\begin{equation}\label{sp:covariance-resolvent}
 S_h=t\int\frac{vv^T}{1+t|v\cdot h|^2}\,d\mu(x,v).
\end{equation}
Let \(H=L^2(\mu;\mathbb C)\). On this space define
\[
 (\mathsf Kf)(x,v)=\int\operatorname{sgn}(y-x)f(y,w)\,d\mu(y,w),
 \qquad \mathsf Qf=\int qf\,d\mu\in\mathbb C^r,
\]
and set
\[
 \mathsf D=M_{\sqrt{1-p}},\qquad
 \mathsf T:H\longrightarrow\mathbb C^r\oplus H,
 \qquad \mathsf Tf=(\mathsf Qf,\mathsf Df).
\]
For an operator between these complex Hilbert spaces, write
\(\mathsf T^{\mathrm{tr}}=J_H\mathsf T^*J_{\mathbb C^r\oplus H}\),
where \(J\) is coordinatewise conjugation. This transpose has the same
operator norm as \(\mathsf T\).

\begin{lemma}[Exact factorization]\label{sp:factorization}
The sign-kernel operator satisfies \(\|\mathsf K\|_{\mathrm{HS}}^2=1\),
and the full score \eqref{mf:score} is
\begin{equation}\label{sp:factorization-eq}
 \mathcal R(h)=
 \|\mathsf T\mathsf K\mathsf T^{\mathrm{tr}}\|_{\mathrm{HS}}^2.
\end{equation}
If \(S_h\preceq I_r\), then \(\mathsf T^*\mathsf T\preceq I_H\).
\end{lemma}
\begin{proof}
The diagonal set \(x=y\) has \(\mu\otimes\mu\)-measure zero because
the position marginal is Lebesgue. Squaring the sign kernel therefore
gives Hilbert--Schmidt norm squared one.

The old--old block \(\mathsf Q\mathsf K\mathsf Q^{\mathrm{tr}}\)
has entries
\[
 \iint\operatorname{sgn}(y-x)q_a(x,v)q_b(y,w)\,d\mu(x,v)d\mu(y,w).
\]
It is the antisymmetric matrix of the ordered wedge \(W_h\). Its
matrix Hilbert--Schmidt norm squared is \(2\|W_h\|^2\).
Each of the two mixed blocks has squared norm
\[
 \int(1-p(v))\|2Q_h(x)-Q_h\|^2\,d\mu(x,v).
\]
The noise--noise block has kernel
\(\sqrt{1-p(v)}\operatorname{sgn}(y-x)\sqrt{1-p(w)}\), so its
squared norm is \((1-m)^2\). Adding the four block norms proves
\eqref{sp:factorization-eq}, including the factor two in the mixed term.

For the contraction assertion, put \(b=q/\sqrt p\) on \(\{p>0\}\)
and \(b=0\) elsewhere. Since \(q=0\) on \(\{p=0\}\),
\[
 \int bb^*\,d\mu
 =t\int_{\{p>0\}}(1-p)vv^T\,d\mu\preceq S_h.
\]
This is an inequality: rows with zero projection on \(h\) can still
contribute to \(S_h\). Weighted Cauchy--Schwarz consequently gives
\[
 \|\mathsf Qf\|^2
 =\left\|\int b\sqrt p\,f\,d\mu\right\|^2
 \le\|S_h\|\int p|f|^2\,d\mu
 \le\int p|f|^2\,d\mu.
\]
Adding \(\|\mathsf Df\|^2=\int(1-p)|f|^2\,d\mu\) proves the claim.
\end{proof}

\begin{lemma}[Quantitative contraction loss]\label{sp:gap}
If \(S_h\preceq I_r\), then
\begin{equation}\label{sp:gap-eq}
 \mathcal R(h)\le1-\frac{m_h^2}{16r}.
\end{equation}
\end{lemma}
\begin{proof}
Put
\(\mathsf P=I_H-\mathsf T^*\mathsf T=M_p-\mathsf Q^*\mathsf Q\succeq0\).
The Hilbert--Schmidt ideal inequality and
Lemma~\ref{sp:factorization} imply
\begin{equation}\label{sp:trace-loss}
 \mathcal R(h)
 \le\|\mathsf T\mathsf K\|_{\mathrm{HS}}^2
 =1-\operatorname{Tr}(\mathsf K^*\mathsf P\mathsf K).
\end{equation}
The conclusion follows immediately when \(m=0\), so assume \(m>0\).
Choose an integer \(d>r\) and write \(w=m/d\). The cumulative
\(p\)-mass is continuous, and \(0\le p\le1\). Hence \([0,1]\)
can be partitioned into \(d\) consecutive intervals of \(p\)-mass
exactly \(w\), each of length at least \(w\).

On interval \(j\), let \(g_j\) rise linearly from zero to one over
its first segment of length \(w/4\), remain one, and fall linearly to
zero over its last segment of length \(w/4\). Extend it by zero and
put \(f_j=-g'_j/2\). Values at the finitely many corners do not matter.
Integration and the Lebesgue position marginal give
\begin{equation}\label{sp:hats}
 \mathsf Kf_j=g_j,\qquad
 \|f_j\|^2=\frac2w,\qquad
 \int p|g_j|^2\,d\mu\ge\frac w2.
\end{equation}
For the last bound, discard the two edge strips. Their total length
is \(w/2\), so they contain at most \(w/2\) of \(p\)-mass; the
remaining mass lies on the plateau.

The disjoint supports make the \(f_j\) orthogonal. Their span \(E\)
has dimension \(d\), and \eqref{sp:hats} gives, for every \(f\in E\),
\[
 \int p|\mathsf Kf|^2\,d\mu\ge\frac{w^2}{4}\|f\|^2.
\]
The subspace \(E\cap\ker(\mathsf Q\mathsf K)\) has dimension at
least \(d-r\). Each unit vector in that subspace satisfies
\[
 \langle\mathsf K^*\mathsf P\mathsf Kf,f\rangle
 =\int p|\mathsf Kf|^2\,d\mu-\|\mathsf Q\mathsf Kf\|^2
 \ge\frac{w^2}{4}.
\]
Summing over \(d-r\) orthonormal vectors and using positivity of the
trace-class operator gives
\[
 \operatorname{Tr}(\mathsf K^*\mathsf P\mathsf K)
 \ge\frac{(d-r)m^2}{4d^2}.
\]
Take \(d=2r\) in \eqref{sp:trace-loss} to obtain
\eqref{sp:gap-eq}.
\end{proof}

\subsection{Stationarity and the third covariance eigenvalue}

\begin{lemma}[An inertia bound]\label{sp:inertia}
Suppose \(h\) is a stationary point of \(\Phi\) and
\(\operatorname{Re}h,\operatorname{Im}h\) are linearly independent.
If \(S_h\) has \(u\) eigenvalues strictly greater than one, counted
with multiplicity, then \(t\Sigma\) has at least \(u+2\) eigenvalues
strictly greater than one.
\end{lemma}
\begin{proof}
Differentiation gives \(S_hh=h\). Since \(S_h\) is real symmetric,
it fixes the two-dimensional real space
\(E=\operatorname{span}_{\mathbb R}(\operatorname{Re}h,
\operatorname{Im}h)\) pointwise. Let \(U\) be the sum of its
eigenspaces with eigenvalue greater than one. Then
\(\dim U=u\) and \(U\perp E\).

The covariance correction
\begin{equation}\label{sp:correction}
 C=t\Sigma-S_h=t\int p_h(v)vv^T\,d\mu\succeq0
\end{equation}
is strictly positive on every nonzero vector \(x\in E\). Indeed,
\(x^TS_hx=\|x\|^2>0\) implies positive measure for
\(\{v:v\cdot x\ne0\}\). On this set \(v\cdot h\ne0\), since
\(x\) is a real linear combination of the real and imaginary parts
of \(h\). Thus \(p_h(v)>0\) there and \(x^TCx>0\).

For a nonzero \(x\in E\oplus U\), a nonzero \(U\)-component gives
\(x^T(S_h-I)x>0\), while \(x^TCx\ge0\). If that component
vanishes, \(x\in E\setminus\{0\}\) and the strict positivity
comes from \(C\). Hence \(t\Sigma-I\) is positive definite on
the \((u+2)\)-dimensional space \(E\oplus U\). The min--max
principle proves the claimed strict eigenvalue count.
\end{proof}

\begin{lemma}\label{sp:maximizer-contraction}
If \(t\lambda_3(\Sigma)\le1\), then \(S_h\preceq I_r\) at
every global maximizer \(h\) of \(\Phi\).
\end{lemma}
\begin{proof}
If \(h\) is real up to a global phase, rotate that phase to make it
real. For every real vector \(z\), the imaginary variation gives
\[
 \frac12\left.\frac{d^2}{d\epsilon^2}
       \Phi(h+i\epsilon z)\right|_{\epsilon=0}
 =z^T(S_h-I)z\le0.
\]
Therefore \(S_h\preceq I\); this also covers \(h=0\).
Otherwise the real and imaginary parts are independent. An eigenvalue
of \(S_h\) above one would, by Lemma~\ref{sp:inertia}, force
\(t\lambda_3(\Sigma)>1\), a contradiction.
\end{proof}

\begin{proof}[Proof of Theorem~\ref{sp:main}]
First suppose an ordered empirical limit exists. At every maximizer,
Lemma~\ref{sp:maximizer-contraction} gives \(S_h\preceq I\).
For a real unit top eigenvector \(z\) of \(\Sigma\),
\[
 1\ge z^TS_hz
 =t\lambda_1(\Sigma)-t\int p_h(v)(v\cdot z)^2\,d\mu
 \ge t\lambda_1(\Sigma)-tL^2m_h.
\]
Consequently
\(m_h\ge(t\lambda_1(\Sigma)-1)/(tL^2)>0\), uniformly over the
maximizing set. Lemma~\ref{sp:gap}, Theorem~\ref{mf:criterion}, and
the endpoint identity \eqref{mf:ratio} now give
\eqref{sp:positive-margin}.

For arbitrary changing orders, every subsequence has a further
subsequence along which the ordered empirical measures converge
weakly on the compact set \([0,1]\times\overline B(0,L)\).
The limit has Lebesgue position marginal and covariance \(\Sigma\).
The lower bound just proved is the same for every such limit, because
it depends only on \(t,L,r,\lambda_1(\Sigma)\). A subsequence
violating \eqref{sp:positive-margin} would contradict its own
ordered-limit bound.
\end{proof}

\begin{corollary}[A necessary three-eigenvalue condition]
\label{sp:negative-necessity}
Under the bounded-row and covariance-convergence assumptions, if
\[
 \liminf_{N\to\infty}
 \frac{D_1(A_N)}{N(N-1)\per A_N}<0,
\]
then \(t\lambda_3(\Sigma)>1\). Thus a negative leading normalized
endpoint derivative requires at least three limiting nontrivial
eigenvalues of \(A_N\) strictly above two.
\end{corollary}
\begin{proof}
If \(t\lambda_1(\Sigma)>1\), this is the contrapositive of
Theorem~\ref{sp:main}. If \(t\lambda_1(\Sigma)\le1\), then
\(\log(1+s)\le s\), with strict inequality for \(s>0\), shows
\(\Phi(h)<0\) for every \(h\ne0\). If the row projection vanishes
almost everywhere, the same conclusion follows from
\(\Phi(h)=-\|h\|^2\). Thus the maximizing set is \(\{0\}\) and
\(\mathcal R(0)=1\). Theorem~\ref{mf:criterion} gives convergence
of the displayed normalized derivative to zero, including the first
critical boundary \(t\lambda_1(\Sigma)=1\). Compactness of ordered
empirical subsequences gives the same conclusion for changing orders.
\end{proof}

\begin{corollary}[Latent rank at most two]\label{sp:rank-two}
For real rows in \(\mathbb R^2\), if \(t\lambda_1(\Sigma)>1\),
then every sequence of orders satisfies
\[
 \liminf_{N\to\infty}
 \frac{D_1(A_N)}{N(N-1)\per A_N}
 \ge\frac{\bigl(t\lambda_1(\Sigma)-1\bigr)^2}{128t^2L^4}>0.
\]
Together with Theorem~\ref{mf:subcritical} and
Corollary~\ref{mf:arbitrary-orders}, this proves eventual endpoint
increase in both noncritical regimes when \(\Sigma\ne0\).
\end{corollary}

\subsection{Scope and the rank-two literature}\label{sp:prior}

The first critical boundary can still leave derivatives of smaller
order than \(N^2\per A_N\); the sign of such derivatives is not
decided here. A public manuscript of liqsweep
\cite{Liqsweep2026} claims permanent-on-top for all real PSD matrices
of rank at most two. If that claim is accepted, its hereditary form,
closure under nonnegative diagonal addition, and the endpoint
sign-Gram identity already imply finite-order endpoint nonnegativity
for that class plus a nonnegative diagonal. The proofs above use no
result from that manuscript. They give an explicit strict asymptotic
margin and apply in arbitrary latent dimension through the number of
supercritical covariance eigenvalues.

% Input-ready conventional-proof component. Labels use the prefix ro:.
% Requires mean_field.tex and the bibliography entry Azuma1967.

\section{Exponential rarity of negative endpoint orders}
\label{ro:section}

The product-limit calculation in Corollary~\ref{mf:product-limit} has
a probabilistic consequence for uniformly permuted deterministic row
clouds. It uses standard bounded-difference concentration in addition
to the full mean-field criterion.

\begin{corollary}[Uniform random orders]\label{ro:random-order}
Fix \(t>0\), \(r\ge1\), and deterministic rows
\(v_{1,N},\ldots,v_{N,N}\in\overline B(0,L)\). Suppose
\[
 \nu_N=\frac1N\sum_{j=1}^N\delta_{v_{j,N}}\Longrightarrow\nu,
 \qquad
 \Sigma=\int vv^T\,d\nu(v),\qquad t\lambda_{\max}(\Sigma)>1.
\]
Let \(\pi_N\) be a uniformly distributed permutation of the \(N\)
row labels, and set
\[
 A_N^{\pi_N}=I_N+(t/N)V_N^{\pi_N}(V_N^{\pi_N})^T.
\]
There exist constants \(\eta,c,C>0\) and
\(N_0\) such that, for \(N\ge N_0\),
\begin{equation}\label{ro:exponential}
 \mathbb P\left\{
 \frac{D_1(A_N^{\pi_N})}{N(N-1)\per A_N^{\pi_N}}<\eta
 \right\}\le Ce^{-cN}.
\end{equation}
In particular the fraction of labelled permutations with nonpositive
endpoint derivative is at most \(Ce^{-cN}\).
Under any joint coupling of the uniform permutations across orders,
the displayed normalized derivative is at least \(\eta\) almost
surely for all sufficiently large \(N\). Independence across orders
is not required.
\end{corollary}
\begin{proof}
Set \(\rho=dx\,\nu(dv)\). For this product ordered limit,
Corollary~\ref{mf:product-limit} gives the strictly positive number
\[
 c_* =\min_{h\in K_\rho}\frac{m_h(4-m_h)}{12}>0.
\]
Take \(\eta=c_*/2\). There is a weak neighborhood \(U\) of
\(\rho\) and an order \(N_1\) such that every bounded-row array
of order \(N\ge N_1\) with ordered empirical measure in \(U\)
has normalized derivative at least \(\eta\). Otherwise, choose
successively smaller neighborhoods of \(\rho\) and increasing
orders to obtain empirical measures converging to \(\rho\), each
with normalized derivative below \(c_*/2\). This contradicts the
liminf bound \(c_*\) from Theorem~\ref{mf:criterion} and
Corollary~\ref{mf:product-limit}.

By the definition of the weak topology on the compact row-position
space, choose finitely many real continuous functions
\(f_1,\ldots,f_J\), \(\epsilon>0\), and
\(B=\max(1,\max_j\|f_j\|_\infty)\) such that
\[
 \left\{\mu:
  \left|\int f_j\,d\mu-\int f_j\,d\rho\right|<\epsilon
                 \text{ for }1\le j\le J\right\}\subseteq U.
\]
For a uniform permutation \(\pi\), define
\[
 X_{j,N}=\frac1N\sum_{i=1}^N f_j(i/N,v_{\pi(i),N}).
\]
Its expectation is
\[
 \mathbb EX_{j,N}
 =\frac1{N^2}\sum_{i,a=1}^N f_j(i/N,v_{a,N})
 \longrightarrow\int f_j\,d\rho.
\]
Indeed the deterministic position-grid measures converge to Lebesgue
measure and \(\nu_N\Longrightarrow\nu\), so their product measures
converge weakly. For large enough \(N\), every expectation above is
within \(\epsilon/2\) of its limit.

Swapping two assigned row labels changes \(X_{j,N}\) by at most
\(4B/N\). Reveal the permutation one position at a time. Conditional
on the already exposed labels, completions for any two choices of
the next label can be coupled by exchanging just those two labels.
Thus the conditional expectations for the two choices differ by at
most \(4B/N\), and each increment of the associated Doob martingale
has absolute value at most \(4B/N\).

The bounded-increment martingale inequality of Azuma
\cite{Azuma1967} now gives
\begin{equation}\label{ro:test-concentration}
 \mathbb P\{|X_{j,N}-\mathbb EX_{j,N}|\ge\epsilon/2\}
 \le2\exp\left(-\frac{\epsilon^2N}{128B^2}\right).
\end{equation}
For clarity, this version follows directly by applying conditional
convexity to a centered increment \(Z\) with \(|Z|\le a\):
\(\mathbb E(e^{sZ}\mid\mathcal F)\le\cosh(sa)
\le e^{s^2a^2/2}\). Iteration and Markov's inequality give the
two-sided bound \(2e^{-u^2/(2Na^2)}\); substitute
\(a=4B/N\) and \(u=\epsilon/2\).

The event in \eqref{ro:exponential} forces the ordered empirical
measure outside \(U\). The union bound in
\eqref{ro:test-concentration} proves the assertion with
\(C=2J\), \(c=\epsilon^2/(128B^2)\), and a sufficiently large
\(N_0\). Finally \(\sum_NCe^{-cN}<\infty\), so the first
Borel--Cantelli lemma gives the almost-sure conclusion under any
joint coupling.
\end{proof}

The constants and the sufficient order in this corollary are
existential: the weak-neighborhood argument does not supply explicit
finite rates for the mean-field approximation. When a deterministic
cloud admits an ordered endpoint counterexample, the corollary shows
that such negative endpoint orders occupy an exponentially small
fraction of its labelled permutations. It makes no assertion of
monotonicity over the full \(q\)-interval for the other orders.


\section{Formal proof and reproducibility}\label{sec:lean}
\subsection{What the Lean development proves}

The accompanying Lean project pins Lean~4.19.0 and Mathlib revision
\begin{center}\small\nolinkurl{c44e0c8ee63ca166450922a373c7409c5d26b00b}.\end{center}
It has $32$ proof modules. A separate audit module enumerates
all $2068$ declarations whose defining module is one of them, including
private and generated declarations, and computes their transitive axioms.
Of these, $1319$ are auxiliary declarations that the Lean~4.19 code
generator stores for executable definitions; beyond the standard axioms
they may refer only to its placeholder \texttt{lcProof} and to one another,
and no other declaration depends on them through an axiom. Each of the other $749$
declarations is a theorem, definition or inductive construction, not an
axiom, and depends only on
\[
 \{\texttt{propext},\ \texttt{Classical.choice},\ \texttt{Quot.sound}\}.
\]
There is no \texttt{sorry}, no added axiom and no native evaluation;
in particular \texttt{Lean.ofReduceBool} does not occur.

The development defines the inversion-weighted polynomial $P_q$ of an
ordered matrix by its permutation sum, exactly as in
Section~\ref{sec:endpoint}, and proves that $D_q$ is its derivative.
It proves, for every real matrix $V$ with three columns, the identity
$2D_1(VV^T)=\binom n2\|f\|_{\F}^2-\|G\|_{\F}^2$ of
Lemma~\ref{lem:endpoint}, with both Fischer norms written as
factorial-weighted coefficient sums of the actual product and wedge
polynomials. It proves the four-row recurrence of
Lemma~\ref{lem:recurrence}, the parity decomposition
\eqref{eq:certP}--\eqref{eq:certT}, the semantics of an executable
sparse integer implementation of the recurrence, and certificate
soundness: for every list $L$ of integer types, if the integer test
$T_U>\binom n2P$ computed from $L$ succeeds, then the Gram matrix of the
row word of $L$ has $D_1<0$, and there are a rational matrix of the same
order that is positive definite and rationals $0<q_1<q_2<1$ with
$P_{q_2}<P_{q_1}$. The last step is a general reduction: a real
positive-semidefinite matrix with negative endpoint derivative yields a
rational positive-definite matrix and a rational interior decreasing
pair. Identity padding is proved for the actual polynomial, together
with preservation of positive definiteness and rationality; one fixed
rational pair therefore serves every larger order.

For the frozen input, the development defines the $\WitnessN$ ordered rows
from the reversed type list, proves that the Gram matrix is positive
semidefinite, and proves that the row matrix and the Gram matrix have
rank exactly three. The rank proof uses the rows with zero-based indices
$0,1,894$, whose matrix is
\[
 \begin{pmatrix}-3&2&20\\-3&-2&20\\3&2&20\end{pmatrix},
 \qquad \det=480,
\]
through a lemma about the outermost type that avoids expanding all
$224$ nested row constructors. This is a different minor from the one
in Table~\ref{tab:witness}; both certify the same rank. Finally it
proves that each of the following hypotheses implies the all-orders
statement
\[
 \exists q_1,q_2\in\mathbb Q,\ 0<q_1<q_2<1,\quad
 \forall N\ge\WitnessN\ \exists A_N\in\mathbb Q^{N\times N},\quad
 A_N\succ0,\quad P_{q_2}(A_N)<P_{q_1}(A_N):
\]
the integer test above for the frozen type list, and the Fischer-norm
inequality $\binom{\WitnessN}2\|f\|_{\F}^2<\|G\|_{\F}^2$ for the frozen rows.

The project also contains scalar lemmas used by the asymptotic
sections: bounds for the spherical Gini gain, logarithmic and radial
errors, rational score margins and condition-number arithmetic for the
families of Theorems~\ref{mf:20002}, \ref{mf:257} and~\ref{rf:228}, and
two finite matrix steps of Section~\ref{sp:section}. These lemmas
formalize the indicated inequalities only; they do not formalize the
probabilistic constructions or limits that supply their hypotheses.

\subsection{What is checked outside Lean}

The evaluation of the integer test on the frozen input is not carried
out in Lean. Kernel reduction of the $224$-step recurrence, whose
polynomials have up to about $10^5$ coefficients of several thousand
digits, is far beyond practical limits. An earlier stage of this work
evaluated it with Lean's native code generator through
\texttt{native\_decide}, which adds \texttt{Lean.ofReduceBool} and hence
the compiler and runtime to the trusted base \cite{LeanNativeDocumentation};
that module is not part of the present development. Instead, the
hypothesis of the conditional theorems is checked by the two exact
programs of Section~\ref{sec:certificate}: the integer implementation of
the four-row recurrence, and the independent row-by-row modular
computation, which determines $P$, $T_U$ and $T$ through the Chinese
remainder theorem. The finite counterexample is therefore a formal
reduction together with an exact computation; it is not a closed
Lean theorem.

The explicit matrix $A=MB+I_n$ and the explicit pair $q_1=1-2/Q$,
$q_2=1-1/Q$ of Section~\ref{fr:section} have conventional proofs; the
formal positive-definite conclusion is existential. The mean-field
endpoint functional, Gaussian concentration, subcritical asymptotic,
spherical saddle analysis, algebraic grids, operator factorization,
spectral obstruction, random-order theorem and coherent-ordering
theorem are conventional proofs as well. Their finite matrix and
scalar ingredients are formalized only where stated above.


\section{Limits and further questions}

The finite certificate provides an economical explicit real witness,
but it does not establish the minimum possible order. The formal
rank-three construction does not prove that rank two is impossible
at every parameter; no minimum-rank theorem is asserted here.
Monotonicity on the whole negative half-interval and failure at an
arbitrarily prescribed positive interior parameter are not settled here.

The mean-field result distinguishes a coherent-saddle threshold from
the sign of the endpoint derivative. The general limiting functional
contains three nonnegative exterior contributions, and negativity of
$D_1$ requires their combined score to exceed one. The subcritical
formula shows that this cannot happen below limiting condition number
two within the stated fixed-rank class. The gap between that obstruction
and the explicit upper construction at $227$ remains open.
The stronger spectral obstruction also requires three limiting
nontrivial eigenvalues above two for a strictly negative $N^2$-scale
derivative. It does not resolve the sign of a vanishing normalized
derivative at the first critical threshold, and is not a theorem about
every finite positive-definite matrix satisfying a spectral bound.
The asymptotic theorems do not supply a numerical sufficient order for
their particular well-conditioned families. The finite matrix and the
bounded-condition constructions should not be conflated.

\section{Research process and AI disclosure}\label{sec:disclosure}

This is fully autonomous agentic research. The finite witness, its
exact certificate and four-row recurrence, the regularization, the
asymptotic and spectral theory, the original Lean development (for
Lean~4.19.0), an evaluation of the certificate with native code, and
the literature search were produced by autonomous AI research agents in
earlier sessions. Their records do not name the system used. The audits
and separate implementations in those records were produced within the
same automated workflow.

Opus 5.5 by Anthropic used with Claude Code assessed the package for
this release. It reran the integer recurrence program on the frozen
inputs, wrote and ran the independent row-by-row modular computation
and the Chinese-remainder determination of Section~\ref{sec:certificate},
checked the endpoint identity and the regularization of
Section~\ref{fr:section} line by line, removed the native evaluation
from the formal development, added the frozen-input module with its
rank and conditional all-orders theorems, ran the fresh Lean build and
whole-project axiom audit, and revised this manuscript. The asymptotic
and spectral sections were carried over from the earlier sessions and
were read, but not rederived line by line, in this release.

Dev D. Goyal is responsible for publication. The work has not been
externally peer reviewed; no expert validation is claimed, AI checks are
not human peer review, and the novelty assessment is provisional and
does not establish priority. Section~\ref{sec:lean} states the formal
scope.

\appendix
\section{Complete ordered integer data}\label{app:data}

The following table is read by increasing index. Each triple is a type
$(a,b,c)$, added as the outer four-row operation in
Section~\ref{sec:certificate}. Thus the list is inner-to-outer. In a
list implementation that recursively wraps the first type around the
tail, the input must be reversed; the formal source records this
conversion explicitly.
\begin{longtable}{@{}r@{\quad}l@{\qquad}r@{\quad}l@{\qquad}r@{\quad}l@{\qquad}r@{\quad}l@{}}
\caption{All $224$ four-row types, in inner-to-outer order.}\label{tab:types}\\
\toprule $j$ & $(a,b,c)$ & $j$ & $(a,b,c)$ & $j$ & $(a,b,c)$ & $j$ & $(a,b,c)$\\\midrule
\endfirsthead
\multicolumn{8}{c}{Table \thetable\ continued}\\\toprule $j$ & $(a,b,c)$ & $j$ & $(a,b,c)$ & $j$ & $(a,b,c)$ & $j$ & $(a,b,c)$\\\midrule
\endhead
\midrule\multicolumn{8}{r}{Continued on the next page}\\\endfoot
\bottomrule\endlastfoot
1 & $(0,1,0)$ & 2 & $(0,1,0)$ & 3 & $(0,1,0)$ & 4 & $(0,1,0)$\\
5 & $(0,1,0)$ & 6 & $(0,1,0)$ & 7 & $(0,1,0)$ & 8 & $(0,1,0)$\\
9 & $(0,1,0)$ & 10 & $(0,1,0)$ & 11 & $(0,1,0)$ & 12 & $(0,1,0)$\\
13 & $(0,1,0)$ & 14 & $(0,1,0)$ & 15 & $(1,23,11)$ & 16 & $(1,23,11)$\\
17 & $(1,30,20)$ & 18 & $(19,24,2)$ & 19 & $(19,24,2)$ & 20 & $(31,32,3)$\\
21 & $(31,32,3)$ & 22 & $(11,25,2)$ & 23 & $(11,25,2)$ & 24 & $(12,31,5)$\\
25 & $(15,28,3)$ & 26 & $(15,28,3)$ & 27 & $(18,28,3)$ & 28 & $(18,28,3)$\\
29 & $(29,16,2)$ & 30 & $(29,16,2)$ & 31 & $(20,9,1)$ & 32 & $(30,25,3)$\\
33 & $(19,16,2)$ & 34 & $(32,22,3)$ & 35 & $(32,22,3)$ & 36 & $(3,23,10)$\\
37 & $(3,23,10)$ & 38 & $(6,29,10)$ & 39 & $(14,5,1)$ & 40 & $(14,5,1)$\\
41 & $(31,7,2)$ & 42 & $(31,7,2)$ & 43 & $(18,2,1)$ & 44 & $(28,3,2)$\\
45 & $(29,1,1)$ & 46 & $(2,23,23)$ & 47 & $(1,11,11)$ & 48 & $(2,24,31)$\\
49 & $(2,24,31)$ & 50 & $(7,21,6)$ & 51 & $(7,21,6)$ & 52 & $(13,29,7)$\\
53 & $(4,27,17)$ & 54 & $(4,27,17)$ & 55 & $(3,27,22)$ & 56 & $(3,27,22)$\\
57 & $(9,15,4)$ & 58 & $(9,15,4)$ & 59 & $(23,31,8)$ & 60 & $(7,29,16)$\\
61 & $(7,29,16)$ & 62 & $(6,19,8)$ & 63 & $(6,19,8)$ & 64 & $(5,11,4)$\\
65 & $(5,11,4)$ & 66 & $(10,17,6)$ & 67 & $(30,31,9)$ & 68 & $(30,31,9)$\\
69 & $(17,21,7)$ & 70 & $(17,21,7)$ & 71 & $(25,21,7)$ & 72 & $(31,26,9)$\\
73 & $(16,11,4)$ & 74 & $(6,25,20)$ & 75 & $(6,25,20)$ & 76 & $(8,21,12)$\\
77 & $(8,21,12)$ & 78 & $(26,15,6)$ & 79 & $(26,15,6)$ & 80 & $(19,9,4)$\\
81 & $(19,7,4)$ & 82 & $(19,7,4)$ & 83 & $(29,7,6)$ & 84 & $(29,7,6)$\\
85 & $(16,31,16)$ & 86 & $(16,31,16)$ & 87 & $(4,12,9)$ & 88 & $(14,2,3)$\\
89 & $(14,2,3)$ & 90 & $(24,1,5)$ & 91 & $(24,1,5)$ & 92 & $(4,6,3)$\\
93 & $(19,28,14)$ & 94 & $(19,24,12)$ & 95 & $(20,21,11)$ & 96 & $(20,21,11)$\\
97 & $(28,24,13)$ & 98 & $(28,24,13)$ & 99 & $(6,23,25)$ & 100 & $(6,23,25)$\\
101 & $(10,29,27)$ & 102 & $(20,32,21)$ & 103 & $(20,32,21)$ & 104 & $(14,30,23)$\\
105 & $(14,30,23)$ & 106 & $(32,20,13)$ & 107 & $(32,20,13)$ & 108 & $(4,3,2)$\\
109 & $(31,12,11)$ & 110 & $(31,12,11)$ & 111 & $(27,14,11)$ & 112 & $(29,15,12)$\\
113 & $(1,7,12)$ & 114 & $(1,7,12)$ & 115 & $(5,21,29)$ & 116 & $(30,6,11)$\\
117 & $(30,6,11)$ & 118 & $(10,3,4)$ & 119 & $(10,3,4)$ & 120 & $(28,32,23)$\\
121 & $(25,29,21)$ & 122 & $(21,19,14)$ & 123 & $(20,1,8)$ & 124 & $(20,1,8)$\\
125 & $(31,3,13)$ & 126 & $(31,3,13)$ & 127 & $(15,26,24)$ & 128 & $(15,26,24)$\\
129 & $(13,28,28)$ & 130 & $(24,25,22)$ & 131 & $(24,25,22)$ & 132 & $(9,12,11)$\\
133 & $(9,12,11)$ & 134 & $(17,13,12)$ & 135 & $(17,13,12)$ & 136 & $(31,19,20)$\\
137 & $(25,10,14)$ & 138 & $(25,10,14)$ & 139 & $(27,14,17)$ & 140 & $(25,13,16)$\\
141 & $(8,20,27)$ & 142 & $(9,22,30)$ & 143 & $(12,25,30)$ & 144 & $(16,3,9)$\\
145 & $(16,3,9)$ & 146 & $(24,8,15)$ & 147 & $(27,9,17)$ & 148 & $(32,25,29)$\\
149 & $(23,18,21)$ & 150 & $(23,22,26)$ & 151 & $(17,20,24)$ & 152 & $(17,20,24)$\\
153 & $(15,23,29)$ & 154 & $(12,18,23)$ & 155 & $(1,5,11)$ & 156 & $(3,14,31)$\\
157 & $(6,17,30)$ & 158 & $(13,1,9)$ & 159 & $(13,1,9)$ & 160 & $(18,2,13)$\\
161 & $(18,2,13)$ & 162 & $(11,6,10)$ & 163 & $(11,6,10)$ & 164 & $(22,17,25)$\\
165 & $(8,3,7)$ & 166 & $(8,3,7)$ & 167 & $(16,7,15)$ & 168 & $(16,7,15)$\\
169 & $(9,9,14)$ & 170 & $(16,16,25)$ & 171 & $(17,12,21)$ & 172 & $(13,19,31)$\\
173 & $(7,10,17)$ & 174 & $(9,17,31)$ & 175 & $(7,13,24)$ & 176 & $(20,15,29)$\\
177 & $(21,16,31)$ & 178 & $(9,9,17)$ & 179 & $(31,1,32)$ & 180 & $(30,1,31)$\\
181 & $(29,1,31)$ & 182 & $(28,1,30)$ & 183 & $(15,1,17)$ & 184 & $(15,1,17)$\\
185 & $(18,1,21)$ & 186 & $(19,2,23)$ & 187 & $(19,2,23)$ & 188 & $(12,1,15)$\\
189 & $(12,1,15)$ & 190 & $(3,10,29)$ & 191 & $(3,10,29)$ & 192 & $(5,10,24)$\\
193 & $(13,12,29)$ & 194 & $(12,11,27)$ & 195 & $(3,4,10)$ & 196 & $(7,9,23)$\\
197 & $(9,0,14)$ & 198 & $(16,0,25)$ & 199 & $(8,0,13)$ & 200 & $(13,0,22)$\\
201 & $(17,0,29)$ & 202 & $(9,0,16)$ & 203 & $(5,0,9)$ & 204 & $(8,0,15)$\\
205 & $(10,0,19)$ & 206 & $(1,0,2)$ & 207 & $(12,0,25)$ & 208 & $(9,0,19)$\\
209 & $(14,0,31)$ & 210 & $(4,0,9)$ & 211 & $(11,0,26)$ & 212 & $(5,0,12)$\\
213 & $(11,0,28)$ & 214 & $(3,6,25)$ & 215 & $(4,7,31)$ & 216 & $(5,6,28)$\\
217 & $(5,6,29)$ & 218 & $(3,0,10)$ & 219 & $(9,0,31)$ & 220 & $(4,0,15)$\\
221 & $(6,0,25)$ & 222 & $(5,0,22)$ & 223 & $(3,2,19)$ & 224 & $(3,2,20)$\\
\end{longtable}


\bibliographystyle{plain}
{\small\raggedright\bibliography{references}}
\end{document}
